% Generated mechanically from the frozen hypercovering.tex target.
% Do not edit this generated file; edit the bound locale source instead.

% OK, start here.
%

\setcounter{chapter}{24}
\chapter{超覆盖}



\phantomsection
\label{hypercovering-section-phantom}


\section{引言}
\label{hypercovering-section-introduction}

\noindent
设 $\mathcal{C}$ 为位点，见《位点》定义 \ref{sites-definition-site}。
设 $X$ 为 $\mathcal{C}$ 的对象。
给定 $\mathcal{C}$ 上的阿贝尔层 $\mathcal{F}$，
我们希望计算其上同调群
$$
H^i(X, \mathcal{F}).
$$
按照一般定义（《位点上的上同调》第
\ref{sites-cohomology-section-cohomology-sheaves} 节），
可通过选取内射分辨
$
0 \to \mathcal{F} \to \mathcal{I}^0 \to \mathcal{I}^1 \to \ldots
$
并置
$$
H^i(X, \mathcal{F})
=
H^i(
\Gamma(X, \mathcal{I}^0) \to
\Gamma(X, \mathcal{I}^1) \to
\Gamma(X, \mathcal{I}^2)\to \ldots)
$$
本章的目标是证明，在 $\mathcal{C}$ 有纤维积的情形下，
也可以不选取内射分辨而计算这些上同调群。
为此我们将使用超覆盖。

\medskip\noindent
位点中的超覆盖是覆盖的推广；见
\cite[Expos\'e V, Sec. 7]{SGA4}。给定对象 $X$ 的超覆盖 $K$，
存在一个从 {\v C}ech 到上同调的谱序列，
它用该层在 $K$ 的各分量 $K_n$ 上的上同调来表示阿贝尔层
$\mathcal{F}$ 在 $X$ 上的上同调。事实表明，超覆盖总是足够多，
因而对所有超覆盖取余极限时，谱序列退化，且 $\mathcal{F}$
在 $X$ 上的上同调由 {\v C}ech 上同调群的余极限计算。

\medskip\noindent
还可考虑一种更一般的构造，即具有上同调下降的单纯增广；见
\cite[Expos\'e Vbis]{SGA4}。关于上同调下降，一份很好的手稿是
Brian Conrad 的文章，见
\url{https://math.stanford.edu/~conrad/papers/hypercover.pdf}。
我们将在关于单纯空间的章节中回到这些问题，并在那里证明，例如，
``局部紧''拓扑空间的真超覆盖具有上同调下降（《单纯空间》第
\ref{spaces-simplicial-section-proper-hypercovering} 节）。
我们的处理方法是把这一陈述归结为本章构造的从 {\v C}ech 到上同调的谱序列。






























\section{半可表对象}
\label{hypercovering-section-semi-representable}

\noindent
首先作如下定义。
字母 ``SR'' 表示半可表（Semi-Representable）。

\begin{definition}
\label{hypercovering-definition-SR}
设 $\mathcal{C}$ 为范畴。用 $\text{SR}(\mathcal{C})$
表示如下定义的{\it 半可表对象}范畴：
\begin{enumerate}
\item 对象是对象族 $\{U_i\}_{i \in I}$；
\item 态射 $\{U_i\}_{i \in I} \to \{V_j\}_{j \in J}$ 由映射
$\alpha : I \to J$ 以及对每个 $i \in I$ 给定的 $\mathcal{C}$ 中态射
$f_i : U_i \to V_{\alpha(i)}$ 构成。
\end{enumerate}
设 $X \in \Ob(\mathcal{C})$ 为 $\mathcal{C}$ 的对象。
{\it $X$ 上的半可表对象}范畴是范畴
$\text{SR}(\mathcal{C}, X) = \text{SR}(\mathcal{C}/X)$.
\end{definition}

\noindent
该定义实质上等价于
\cite[Expos\'e V, Subsection 7.3.0]{SGA4}。注意，这是一个``大''范畴。
稍后我们将``限制''构造 $X$ 的超覆盖所需指标集 $I$ 的大小，
届时便可重新定义 $\text{SR}(\mathcal{C}, X)$，使之成为一个范畴。
现明确写出 $\text{SR}(\mathcal{C}, X)$ 的对象与态射：
\begin{enumerate}
\item 对象是态射族 $\{U_i \to X\}_{i \in I}$；
\item 态射 $\{U_i \to X\}_{i \in I} \to
\{V_j \to X\}_{j \in J}$ 由映射 $\alpha : I \to J$
以及对每个 $i \in I$ 给定的 $X$ 上态射
$f_i : U_i \to V_{\alpha(i)}$ 构成。
\end{enumerate}
存在遗忘函子
$\text{SR}(\mathcal{C}, X) \to \text{SR}(\mathcal{C})$.

\begin{definition}
\label{hypercovering-definition-SR-F}
设 $\mathcal{C}$ 为范畴。
用 $F$ 表示{\it 把预层赋给半可表对象的函子}。以公式表示为
\begin{eqnarray*}
F : \text{SR}(\mathcal{C}) & \longrightarrow & \textit{PSh}(\mathcal{C}) \\
\{U_i\}_{i \in I} & \longmapsto & \amalg_{i\in I} h_{U_i}
\end{eqnarray*}
其中 $h_U$ 表示与对象 $U$ 相伴的可表预层。
\end{definition}

\noindent
给定态射 $U \to X$，得到可表预层的态射 $h_U \to h_X$。
因此，我们常把 $F$ 在 $\text{SR}(\mathcal{C}, X)$ 上的限制视为取值于
$h_X$ 上集合预层范畴、即 $\textit{PSh}(\mathcal{C})/h_X$ 的函子。
相应图示为：
$$
\xymatrix{
\text{SR}(\mathcal{C}, X) \ar[r]_F \ar[d] &
\textit{PSh}(\mathcal{C})/h_X \ar[d] \\
\text{SR}(\mathcal{C}) \ar[r]^F &
\textit{PSh}(\mathcal{C})
}
$$
下面讨论半可表对象范畴中极限的存在性。

\begin{lemma}
\label{hypercovering-lemma-coprod-prod-SR}
设 $\mathcal{C}$ 为范畴。
\begin{enumerate}
\item 范畴 $\text{SR}(\mathcal{C})$ 有余积，且 $F$ 与余积交换；
\item 函子 $F : \text{SR}(\mathcal{C}) \to \textit{PSh}(\mathcal{C})$
与极限交换；
\item 若 $\mathcal{C}$ 有纤维积，则 $\text{SR}(\mathcal{C})$ 有纤维积；
\item 若 $\mathcal{C}$ 有二元积，则 $\text{SR}(\mathcal{C})$ 有二元积；
\item 若 $\mathcal{C}$ 有等化子，则 $\text{SR}(\mathcal{C})$ 也有等化子；
\item 若 $\mathcal{C}$ 有终对象，则 $\text{SR}(\mathcal{C})$ 也有终对象。
\end{enumerate}
设 $X \in \Ob(\mathcal{C})$。
\begin{enumerate}
\item 范畴 $\text{SR}(\mathcal{C}, X)$ 有余积，且 $F$ 与余积交换；
\item 若 $\mathcal{C}$ 有纤维积，则 $\text{SR}(\mathcal{C}, X)$
有有限极限，且
$F : \text{SR}(\mathcal{C}, X) \to \textit{PSh}(\mathcal{C})/h_X$
与这些极限交换。
\end{enumerate}
\end{lemma}

\begin{proof}
先证明关于 $\text{SR}(\mathcal{C})$ 的各项结论。
证明 (1)。$\{U_i\}_{i \in I}$ 与 $\{V_j\}_{j \in J}$ 的余积是
$\{U_i\}_{i \in I} \amalg \{V_j\}_{j \in J}$；换言之，它是以
$I \amalg J$ 为指标集的对象族，对元素 $k \in I \amalg J$，若
$k = i \in I$ 则给出 $U_i$，若 $k = j \in J$ 则给出 $V_j$。
对象族之族的余积也类似。显然 $F$ 与这些余积交换。

\medskip\noindent
证明 (2)。对 $\Ob(\mathcal{C})$ 中的 $U$，考虑
$\text{SR}(\mathcal{C})$ 的对象 $\{U\}$。显然
$\Mor_{\text{SR}(\mathcal{C})}(\{U\}, K)) = F(K)(U)$
对 $K \in \Ob(\text{SR}(\mathcal{C}))$ 成立。由于预层的极限在截面层面计算
（《位点》第 \ref{sites-section-limits-colimits-PSh} 节），
故 $F$ 与极限交换。

\medskip\noindent
证明 (3)。设给定态射
$(\alpha, f_i) : \{U_i\}_{i \in I} \to \{V_j\}_{j \in J}$
以及态射
$(\beta, g_k) : \{W_k\}_{k \in K} \to \{V_j\}_{j \in J}$.
这些态射的纤维积由下式给出：
$$
\{ U_i \times_{f_i, V_j, g_k} W_k\}_{(i, j, k) \in I \times J \times K
\text{使得 } j = \alpha(i) = \beta(k)}
$$
若 $\mathcal{C}$ 有纤维积，则这些纤维积存在。

\medskip\noindent
证明 (4)。$\{U_i\}_{i \in I}$ 与 $\{V_j\}_{j \in J}$ 的积为
$\{U_i \times V_j\}_{i \in I, j \in J}$。若 $\mathcal{C}$ 有积，
则这些积存在。

\medskip\noindent
证明 (5)。两个映射
$(\alpha, f_i), (\alpha', f'_i) : \{U_i\}_{i \in I} \to \{V_j\}_{j \in J}$
的等化子为
$$
\{
\text{Eq}(f_i, f'_i : U_i \to V_{\alpha(i)})
\}_{i \in I,\ \alpha(i) = \alpha'(i)}
$$
若 $\mathcal{C}$ 有等化子，则这些等化子存在。

\medskip\noindent
证明 (6)。若 $X$ 是 $\mathcal{C}$ 的终对象，则 $\{X\}$ 是
$\text{SR}(\mathcal{C})$ 的终对象。

\medskip\noindent
下面证明关于 $\text{SR}(\mathcal{C}, X)$ 的陈述。
把上述结果应用于范畴 $\mathcal{C}/X$，并利用
$\text{SR}(\mathcal{C}/X) = \text{SR}(\mathcal{C}, X)$ 以及
$\textit{PSh}(\mathcal{C}/X) = \textit{PSh}(\mathcal{C})/h_X$
（把《位点》引理 \ref{sites-lemma-essential-image-j-shriek}
应用于赋予混沌拓扑的 $\mathcal{C}$），即可得到这些陈述。
不过我们也直接论证如下。显然，
$\{U_i \to X\}_{i \in I}$ 与 $\{V_j \to X\}_{j \in J}$ 的余积是
$\{U_i \to X\}_{i \in I} \amalg \{V_j \to X\}_{j \in J}$；
终点为 $X$ 的态射族之族的余积也类似。对象 $\{X \to X\}$ 是
$\text{SR}(\mathcal{C}, X)$ 的终对象。设给定态射
$(\alpha, f_i) : \{U_i \to X\}_{i \in I} \to \{V_j \to X\}_{j \in J}$
以及态射
$(\beta, g_k) : \{W_k \to X\}_{k \in K} \to \{V_j \to X\}_{j \in J}$.
这些态射的纤维积由下式给出：
$$
\{ U_i \times_{f_i, V_j, g_k} W_k \to X \}_{(i, j, k) \in I \times J \times K
\text{使得 } j = \alpha(i) = \beta(k)}
$$
由假设 $\mathcal{C}$ 有纤维积，这些纤维积存在。
因此 $\text{SR}(\mathcal{C}, X)$ 有有限极限；
见《范畴》引理 \ref{categories-lemma-finite-limits-exist}。
略去此情形下关于函子 $F$ 的陈述的验证。
\end{proof}





\section{超覆盖}
\label{hypercovering-section-hypercoverings}

\noindent
若再假设我们的范畴是位点，便可作如下定义。

\begin{definition}
\label{hypercovering-definition-covering-SR}
设 $\mathcal{C}$ 为位点。设
$f = (\alpha, f_i) : \{U_i\}_{i \in I} \to \{V_j\}_{j \in J}$
为范畴 $\text{SR}(\mathcal{C})$ 中的态射。若对每个 $j \in J$，
态射族 $\{U_i \to V_j\}_{i \in I, \alpha(i) = j}$
都是位点 $\mathcal{C}$ 的覆盖，则称 $f$ 为{\it 覆盖}。
设 $X$ 为 $\mathcal{C}$ 的对象。若 $\text{SR}(\mathcal{C}, X)$
中的态射 $K \to L$ 在 $\text{SR}(\mathcal{C})$ 中的像是覆盖，
则称它为{\it 覆盖}。
\end{definition}

\begin{lemma}
\label{hypercovering-lemma-covering-permanence}
设 $\mathcal{C}$ 为位点。
\begin{enumerate}
\item $\text{SR}(\mathcal{C})$ 中覆盖的复合仍为覆盖。
\item 若 $K \to L$ 是 $\text{SR}(\mathcal{C})$ 中的覆盖，
而 $L' \to L$ 是态射，则 $L' \times_L K$ 存在，且
$L' \times_L K \to L'$ 是覆盖。
\item 若 $\mathcal{C}$ 有二元积，并且 $A \to B$ 与 $K \to L$
是 $\text{SR}(\mathcal{C})$ 中的覆盖，则
$A \times K \to B \times L$ 是覆盖。
\end{enumerate}
设 $X \in \Ob(\mathcal{C})$。则 (1) 与 (2) 对
$\text{SR}(\mathcal{C}, X)$ 成立；若 $\mathcal{C}$ 有纤维积，
则 (3) 也成立。
\end{lemma}

\begin{proof}
(1) 由位点的公理立即得到。(2) 由引理
\ref{hypercovering-lemma-coprod-prod-SR} 的证明中 $\text{SR}(\mathcal{C})$
内纤维积的构造，以及 $\mathcal{C}$ 的覆盖中各态射须为可表态射的要求得到。
把 $A \times K \to B \times L$ 看作复合
$A \times K \to B \times K \to B \times L$，它因而是覆盖的基变换的复合，
故得 (3)。最后一个陈述由
$\text{SR}(\mathcal{C}, X) = \text{SR}(\mathcal{C}/X)$ 得出。
\end{proof}

\noindent
由引理 \ref{hypercovering-lemma-coprod-prod-SR} 与《单纯方法》引理
\ref{simplicial-lemma-existence-cosk}，若 $\mathcal{C}$ 有纤维积，
则 $\text{SR}(\mathcal{C}, X)$ 的截断单纯对象的余骨架存在。
因此下述定义是有意义的。

\begin{definition}
\label{hypercovering-definition-hypercovering}
设 $\mathcal{C}$ 为有纤维积的位点，且 $X \in \Ob(\mathcal{C})$。
{\it $X$ 的超覆盖}是 $\text{SR}(\mathcal{C}, X)$ 的单纯对象 $K$，满足：
\begin{enumerate}
\item 对象 $K_0$ 是位点 $\mathcal{C}$ 中 $X$ 的覆盖；
\item 对每个 $n \geq 0$，典范态射
$$
K_{n + 1} \longrightarrow (\text{cosk}_n \text{sk}_n K)_{n + 1}
$$
是上述意义下的覆盖。
\end{enumerate}
\end{definition}

\noindent
条件 (1) 是有意义的，因为 $\text{SR}(\mathcal{C}, X)$ 的每个对象
毕竟都是终点为 $X$ 的态射族。它也可以表述为：从 $K_0$ 到
$\text{SR}(\mathcal{C}, X)$ 的终对象的态射是覆盖。

\begin{example}[{\v C}ech 超覆盖]
\label{hypercovering-example-cech}
设 $\mathcal{C}$ 为有纤维积的位点，且
$\{U_i \to X\}_{i \in I}$ 是 $\mathcal{C}$ 的覆盖。置
$K_0 = \{U_i \to X\}_{i \in I}$。则 $K_0$ 是
$\text{SR}(\mathcal{C}, X)$ 的 $0$-截断单纯对象，故可构造
$$
K = \text{cosk}_0 K_0.
$$
显然 $K$ 满足定义 \ref{hypercovering-definition-hypercovering} 的条件 (1)。
由《单纯方法》引理 \ref{simplicial-lemma-cosk-up}，所有态射
$K_{n + 1} \to (\text{cosk}_n \text{sk}_n K)_{n + 1}$ 都是同构，
故它也满足条件 (2)。注意，各项 $K_n$ 正是通常的
$$
K_n = \{
U_{i_0} \times_X U_{i_1} \times_X \ldots \times_X U_{i_n} \to X
\}_{(i_0, i_1, \ldots, i_n) \in I^{n + 1}}
$$
这种形式的 $X$ 的超覆盖称为 $X$ 的{\it {\v C}ech 超覆盖}。
\end{example}

\begin{example}[由位点的单纯对象给出的超覆盖]
\label{hypercovering-example-hypercovering-in-C}
设 $\mathcal{C}$ 为有纤维积的位点，$X \in \Ob(\mathcal{C})$，
并设 $U$ 为 $\mathcal{C}$ 的单纯对象。仍记 $U_n = U([n])$。
最后，假设给定增广
$$
a : U \to X
$$
在此情形下，可考虑 $\text{SR}(\mathcal{C}, X)$ 的单纯对象 $K$，
其各项为 $K_n = \{U_n \to X\}$。当且仅当下列三个条件成立时，
$K$ 是定义 \ref{hypercovering-definition-hypercovering} 意义下的 $X$ 的超覆盖
\footnote{由于 $\mathcal{C}$ 有纤维积，范畴 $\mathcal{C}/X$
有全部有限极限。因此，由《单纯方法》引理
\ref{simplicial-lemma-existence-cosk}，所需余骨架存在。}：
\begin{enumerate}
\item $\{U_0 \to X\}$ 是 $\mathcal{C}$ 的覆盖；
\item $\{U_1 \to U_0 \times_X U_0\}$ 是 $\mathcal{C}$ 的覆盖；
\item $\{U_{n + 1} \to (\text{cosk}_n\text{sk}_n U)_{n + 1}\}$
对 $n \geq 1$ 是 $\mathcal{C}$ 的覆盖。
\end{enumerate}
略去直接的验证。
\end{example}

\begin{example}[与一个覆盖相伴的 {\v C}ech 超覆盖]
\label{hypercovering-example-cech-cover}
设 $\mathcal{C}$ 为有纤维积的位点。设 $U \to X$ 是
$\mathcal{C}$ 的态射，使得 $\{U \to X\}$ 是 $\mathcal{C}$ 的覆盖
\footnote{$\mathcal{C}$ 中具有此性质的态射有时称为``覆盖态射''。}。
考虑 $\text{SR}(\mathcal{C}, X)$ 的单纯对象 $K$，其各项为
$$
K_n = \{U \times_X U \times_X \ldots \times_X U \to X\}
\quad (n + 1 \text{ 个因子})
$$
则 $K$ 是 $X$ 的超覆盖。该例同时是例 \ref{hypercovering-example-cech}
与例 \ref{hypercovering-example-hypercovering-in-C} 的特殊情形。
\end{example}

\begin{lemma}
\label{hypercovering-lemma-hypercoverings-set}
设 $\mathcal{C}$ 为有纤维积的位点，且 $X \in \Ob(\mathcal{C})$。
$X$ 的所有超覆盖所成的类是集合。
\end{lemma}

\begin{proof}
由于 $\mathcal{C}$ 是位点，$X$ 的所有覆盖所成的类是集合。
因此，可取的 $K_0$ 所成的类是集合。假设已经证明可取的
$K_0, \ldots, K_n$ 所成的类是集合。于是只需证明，给定
$K_0, \ldots, K_n$ 后，可取的 $K_{n + 1}$ 所成的类是集合。
这显然成立，因为 $K_{n + 1}$ 必须从
$(\text{cosk}_n \text{sk}_n K)_{n + 1}$ 的所有可能覆盖中选取。
\end{proof}

\begin{remark}
\label{hypercovering-remark-hypercoverings-really-set}
该引理不只是说，存在由超覆盖的选择组成且为集合的共尾系统；
它断言超覆盖本身确实组成一个集合。
\end{remark}

\noindent
$\mathcal{C}$ 上的预层范畴有有限（余）极限。因此，
集合预层上存在函子 $\text{cosk}_n$。

\begin{lemma}
\label{hypercovering-lemma-hypercovering-F}
设 $\mathcal{C}$ 为有纤维积的位点，$X \in \Ob(\mathcal{C})$，
并设 $K$ 为 $X$ 的超覆盖。考虑 $\textit{PSh}(\mathcal{C})$
的单纯对象 $F(K)$，并赋予它到常值单纯预层 $h_X$ 的增广。
\begin{enumerate}
\item 预层态射 $F(K)_0 \to h_X$ 在层化后成为满射。
\item 态射
$$
(d^1_0, d^1_1) :
F(K)_1
\longrightarrow
F(K)_0 \times_{h_X} F(K)_0
$$
在层化后成为满射。
\item 对每个 $n \geq 1$，态射
$$
F(K)_{n + 1} \longrightarrow (\text{cosk}_n \text{sk}_n F(K))_{n + 1}
$$
在层化后成为满射。
\end{enumerate}
\end{lemma}

\begin{proof}
我们将使用如下事实：若 $\{U_i \to U\}_{i \in I}$
是位点 $\mathcal{C}$ 的覆盖，则态射
$$
\amalg_{i \in I} h_{U_i} \to h_U
$$
在层化后成为满射；见《位点》引理
\ref{sites-lemma-covering-surjective-after-sheafification}。
故第一个断言立即成立。

\medskip\noindent
对第二个断言，由《单纯方法》例 \ref{simplicial-example-cosk0}，
单纯对象 $\text{cosk}_0 \text{sk}_0 K$ 的各项为
$K_0 \times \ldots \times K_0$。因此，由超覆盖的定义，
$(d^1_0, d^1_1) : K_1 \to K_0 \times K_0$ 是覆盖。
结合上述断言以及 $F$ 把积变为 $h_X$ 上的纤维积这一事实，便得到 (2)。

\medskip\noindent
对第三个断言，我们声称
$\text{cosk}_n \text{sk}_n F(K) =
F(\text{cosk}_n \text{sk}_n K)$
对 $n \geq 1$ 成立。为证明这一点，暂用 $F'$ 表示函子
$\text{SR}(\mathcal{C}, X) \to \textit{PSh}(\mathcal{C})/h_X$.
由引理 \ref{hypercovering-lemma-coprod-prod-SR}，函子 $F'$ 与有限极限交换。
根据《单纯方法》第 \ref{simplicial-section-skeleton} 节中
对 $\text{cosk}_n$ 函子的描述，有
$\text{cosk}_n \text{sk}_n F'(K) =
F'(\text{cosk}_n \text{sk}_n K)$.
回忆《单纯方法》引理 \ref{simplicial-lemma-existence-cosk}
中描述 $(\text{cosk}_n U)_m$ 时所用的范畴是
$(\Delta/[m])^{opp}_{\leq n}$。一个有趣的练习是证明，只要
$n \geq 1$，范畴 $(\Delta/[m])_{\leq n}$ 就是连通的（见《范畴》定义
\ref{categories-definition-category-connected}）。因此，由《范畴》引理
\ref{categories-lemma-connected-limit-over-X}，有
$\text{cosk}_n \text{sk}_n F'(K) =
\text{cosk}_n \text{sk}_n F(K)$，从而该断言成立。
由此得到性质 (2)，因为此时可见，(2) 中的态射是把函子 $F$
应用于定义 \ref{hypercovering-definition-covering-SR} 意义下的覆盖所得；
结论遂由本证明开头提到的第一个事实得出。
\end{proof}



\section{无环性}
\label{hypercovering-section-acyclicity}

\noindent
设 $\mathcal{C}$ 为位点。对集合预层 $\mathcal{F}$，用
$\mathbf{Z}_\mathcal{F}$ 表示由下述规则定义的阿贝尔群预层：
$$
\mathbf{Z}_\mathcal{F}(U) = \text{由 }\mathcal{F}(U)\text{ 生成的自由阿贝尔群}.
$$
有时称它为{\it $\mathcal{F}$ 上的自由阿贝尔预层}。
当然，构造 $\mathcal{F} \mapsto \mathbf{Z}_\mathcal{F}$ 是函子，
且是遗忘函子
$\textit{PAb}(\mathcal{C}) \to \textit{PSh}(\mathcal{C})$.
的左伴随。当然，层化 $\mathbf{Z}_\mathcal{F}^\#$ 是阿贝尔群层，
而函子 $\mathcal{F} \mapsto \mathbf{Z}_\mathcal{F}^\#$ 也为左伴随。
有时称 $\mathbf{Z}_\mathcal{F}^\#$ 为
{\it $\mathcal{F}$ 上的自由阿贝尔层}。

\medskip\noindent
对位点 $\mathcal{C}$ 的对象 $X$，用 $\mathbf{Z}_X$ 表示
$h_X$ 上的自由阿贝尔预层，并用 $\mathbf{Z}_X^\#$ 表示其层化。

\begin{definition}
\label{hypercovering-definition-homology}
设 $\mathcal{C}$ 为位点，$K$ 为 $\textit{PSh}(\mathcal{C})$
的单纯对象。由上可得 $\textit{Ab}(\mathcal{C})$ 的单纯对象
$\mathbf{Z}_K^\#$。可取其相伴的阿贝尔预层复形
$s(\mathbf{Z}_K^\#)$；见《单纯方法》第
\ref{simplicial-section-complexes} 节。$K$ 的{\it 同调}是
阿贝尔层复形 $s(\mathbf{Z}_K^\#)$ 的同调。
\end{definition}

\noindent
换言之，$K$ 的{\it 第 $i$ 个同调 $H_i(K)$}是阿贝尔群层
$H_i(K) = H_i(s(\mathbf{Z}_K^\#))$。本节研究 $K$ 是
$\mathcal{C}$ 的对象 $X$ 的超覆盖时的同调。

\begin{lemma}
\label{hypercovering-lemma-compare-cosk0}
设 $\mathcal{C}$ 为位点，$\mathcal{F} \to \mathcal{G}$
为集合预层的态射。用 $K$ 表示 $\textit{PSh}(\mathcal{C})$
的单纯对象，其第 $n$ 项是 $\mathcal{F}$ 在 $\mathcal{G}$ 上的
$(n + 1)$ 重纤维积；见《单纯方法》例
\ref{simplicial-example-fibre-products-simplicial-object}。
若 $\mathcal{F} \to \mathcal{G}$ 在层化后为满射，则
$$
H_i(K) =
\left\{
\begin{matrix}
0 & \text{若} & i > 0\\
\mathbf{Z}_\mathcal{G}^\# & \text{若} & i = 0
\end{matrix}
\right.
$$
次数 $0$ 上的同构由态射
$H_0(K) \to \mathbf{Z}_\mathcal{G}^\#$ 给出；该态射来自映射
$(\mathbf{Z}_K^\#)_0 =
\mathbf{Z}_\mathcal{F}^\# \to \mathbf{Z}_\mathcal{G}^\#$。
\end{lemma}

\begin{proof}
令 $\mathcal{G}' \subset \mathcal{G}$ 为态射
$\mathcal{F} \to \mathcal{G}$ 的像。取 $U \in \Ob(\mathcal{C})$，置
$A = \mathcal{F}(U)$ 和 $B = \mathcal{G}'(U)$。
则单纯集 $K(U)$ 等于第 $n$ 个单形由下式给出的单纯集：
$$
A \times_B A \times_B \ldots \times_B A\ (n + 1 \text{ 个因子)}.
$$
由《单纯方法》引理 \ref{simplicial-lemma-cosk-minus-one-equivalence}，
态射 $K(U) \to B$ 是平凡 Kan 纤维化，因而是同伦等价
（《单纯方法》引理 \ref{simplicial-lemma-trivial-kan-homotopy}）。
对它应用``生成自由阿贝尔群''函子，得到
$$
\mathbf{Z}_K(U) \longrightarrow \mathbf{Z}_B
$$
是同伦等价。注意，$s(\mathbf{Z}_B)$ 是复形
$$
\ldots \to
\bigoplus\nolimits_{b \in B}\mathbf{Z} \xrightarrow{0}
\bigoplus\nolimits_{b \in B}\mathbf{Z} \xrightarrow{1}
\bigoplus\nolimits_{b \in B}\mathbf{Z} \xrightarrow{0}
\bigoplus\nolimits_{b \in B}\mathbf{Z} \to 0
$$
；见《单纯方法》引理 \ref{simplicial-lemma-homology-eilenberg-maclane}。
因此，对 $i > 0$ 有
$H_i(s(\mathbf{Z}_K(U))) = 0$，且
$H_0(s(\mathbf{Z}_K(U))) = \bigoplus_{b \in B}\mathbf{Z}
= \bigoplus_{s \in \mathcal{G}'(U)} \mathbf{Z}$.
这些等同与限制映射相容。

\medskip\noindent
由此推出 $H_i(s(\mathbf{Z}_K)) = 0$ 对 $i > 0$ 成立，且
$H_0(s(\mathbf{Z}_K)) = \mathbf{Z}_{\mathcal{G}'}$；这里在
$\textit{PAb}(\mathcal{C})$ 中计算同调群。由于层化是正合函子，
便得到引理的结论。确切地说，正合性蕴含
$H_0(s(\mathbf{Z}_K))^\# = H_0(s(\mathbf{Z}_K^\#))$，
其他指标也同样。
\end{proof}

\begin{lemma}
\label{hypercovering-lemma-acyclicity}
设 $\mathcal{C}$ 为位点，$f : L \to K$ 为
$\textit{PSh}(\mathcal{C})$ 的单纯对象之间的态射，
且 $n \geq 0$ 为整数。假设：
\begin{enumerate}
\item 对 $i < n$，态射 $L_i \to K_i$ 是同构；
\item 态射 $L_n \to K_n$ 在层化后为满射；
\item 典范映射 $L \to \text{cosk}_n \text{sk}_n L$ 是同构；
\item 典范映射 $K \to \text{cosk}_n \text{sk}_n K$ 是同构。
\end{enumerate}
则 $H_i(f) : H_i(L) \to H_i(K)$ 是同构。
\end{lemma}

\begin{proof}
该证明与上面引理 \ref{hypercovering-lemma-compare-cosk0} 的证明完全相同。
确切地说，先令 $K_n' \subset K_n$ 为映射 $L_n \to K_n$
之像所成的子预层。假设 (2) 意味着 $K_n'$ 的层化等于 $K_n$ 的层化。
此外，由于对所有 $i < n$ 都有 $L_i = K_i$，取
$U_0 = L_0 = K_0, \ldots, U_{n - 1} = L_{n - 1} = K_{n - 1}, U_n = K'_n$.
便得到 $n$-截断单纯预层 $U$。记单纯预层
$K' = \text{cosk}_n U$。由于可将 $K'_m$ 构造为有限极限，
而层化又正合，故 $(K'_m)^\# = K_m$。换言之，
$(K')^\# = K^\#$。再次由层化的正合性，推出
$H_i(K) = H_i(K')$。因此，只需对态射 $L \to K'$ 证明该引理；
也就是说，可以假设 $L_n \to K_n$ 是{\it 预层}的满射！

\medskip\noindent
在此情形下，对 $\mathcal{C}$ 的任意对象 $U$，单纯集的态射
$$
L(U) \longrightarrow K(U)
$$
满足《单纯方法》引理 \ref{simplicial-lemma-section} 的全部假设。
故它是平凡 Kan 纤维化，尤其是同伦等价
（《单纯方法》引理 \ref{simplicial-lemma-trivial-kan-homotopy}）。因此
$$
\mathbf{Z}_L(U) \longrightarrow \mathbf{Z}_K(U)
$$
也是同伦等价。这对所有 $U$ 都成立，结论随之得出。
\end{proof}

\begin{lemma}
\label{hypercovering-lemma-acyclic-hypercover-sheaves}
设 $\mathcal{C}$ 为位点，$K$ 为单纯预层，$\mathcal{G}$ 为预层，
并设 $K \to \mathcal{G}$ 为 $K$ 到 $\mathcal{G}$ 的增广。假设：
\begin{enumerate}
\item 预层态射 $K_0 \to \mathcal{G}$ 在层化后成为满射；
\item 态射
$$
(d^1_0, d^1_1) :
K_1
\longrightarrow
K_0 \times_\mathcal{G} K_0
$$
在层化后成为满射；
\item 对每个 $n \geq 1$，态射
$$
K_{n + 1} \longrightarrow (\text{cosk}_n \text{sk}_n K)_{n + 1}
$$
在层化后成为满射。
\end{enumerate}
则对 $i > 0$ 有 $H_i(K) = 0$，且
$H_0(K) = \mathbf{Z}_\mathcal{G}^\#$。
\end{lemma}

\begin{proof}
对 $n \geq 1$，记 $K^n = \text{cosk}_n \text{sk}_n K$。
把 $K^0$ 定义为这样的单纯对象：$(K^0)_n$ 等于 $(n + 1)$ 重纤维积
$K_0 \times_\mathcal{G} \ldots \times_\mathcal{G} K_0$；
见《单纯方法》例 \ref{simplicial-example-fibre-products-simplicial-object}。
存在态射
$$
K \longrightarrow \ldots \to K^n \to K^{n - 1} \to \ldots \to K^1 \to K^0.
$$
态射 $K \to K^i$ 以及 $j \geq i \geq 1$ 时的 $K^j \to K^i$
来自各函子 $\text{cosk}_n$ 的泛性质。态射 $K^1 \to K^0$
是《单纯方法》评注 \ref{simplicial-remark-augmentation} 中的典范态射。
还回忆，$K^0 \to \text{cosk}_1 \text{sk}_1 K^0$ 是同构；
见《单纯方法》引理 \ref{simplicial-lemma-cosk-minus-one}。

\medskip\noindent
由引理 \ref{hypercovering-lemma-compare-cosk0}，对 $i > 0$ 有
$H_i(K^0) = 0$，且 $H_0(K^0) = \mathbf{Z}_\mathcal{G}^\#$。

\medskip\noindent
取 $n \geq 1$，考虑态射 $K^n \to K^{n - 1}$。
它在次数 $< n$ 的各项上是同构。注意，
$K^n \to \text{cosk}_n \text{sk}_n K^n$ 与
$K^{n - 1} \to \text{cosk}_n \text{sk}_n K^{n - 1}$ 都是同构；又有
$(K^n)_n = K_n$ 及
$(K^{n - 1})_n = (\text{cosk}_{n - 1} \text{sk}_{n - 1} K)_n$。
因此，由假设，$(K^n)_n \to (K^{n - 1})_n$ 是层化后成为满射的预层态射。
由引理 \ref{hypercovering-lemma-acyclicity}，得到
$H_i(K^n) = H_i(K^{n - 1})$。结合上文即证明该引理。
\end{proof}

\begin{lemma}
\label{hypercovering-lemma-hypercovering-acyclic}
设 $\mathcal{C}$ 为有纤维积的位点，$X$ 为 $\mathcal{C}$ 的对象，
并设 $K$ 为 $X$ 的超覆盖。单纯预层 $F(K)$ 的同调在次数 $> 0$
时为 $0$，在次数 $0$ 时等于 $\mathbf{Z}_X^\#$。
\end{lemma}

\begin{proof}
合用引理 \ref{hypercovering-lemma-acyclic-hypercover-sheaves}
与 \ref{hypercovering-lemma-hypercovering-F}。
\end{proof}












\section{{\v C}ech 上同调与超覆盖}
\label{hypercovering-section-hyper-cech}

\noindent
设 $\mathcal{C}$ 为位点。考虑位点 $\mathcal{C}$ 上的阿贝尔群预层
$\mathcal{F}$。它定义函子
\begin{eqnarray*}
\mathcal{F} : \text{SR}(\mathcal{C})^{opp}
& \longrightarrow &
\textit{Ab} \\
\{U_i\}_{i \in I} &
\longmapsto &
\prod\nolimits_{i \in I} \mathcal{F}(U_i)
\end{eqnarray*}
于是，$\text{SR}(\mathcal{C})$ 的单纯对象 $K$ 被变为
$\textit{Ab}$ 的余单纯对象 $\mathcal{F}(K)$。与 $\mathcal{F}(K)$
相伴的上链复形 $s(\mathcal{F}(K))$（《单纯方法》第
\ref{simplicial-section-dold-kan-cosimplicial} 节）称为
$\mathcal{F}$ 关于单纯对象 $K$ 的 {\v C}ech 复形。置
$$
\check{H}^i(K, \mathcal{F})
=
H^i(s(\mathcal{F}(K))).
$$
并称之为 $\mathcal{F}$ 关于 $K$ 的第 $i$ 个 {\v C}ech 上同调群。
本节证明《层上同调》第 \ref{cohomology-section-cech}、
\ref{cohomology-section-cech-functor} 与
\ref{cohomology-section-cech-cohomology-cohomology} 节中
关于开覆盖的 {\v C}ech 上同调若干结果的类似物。

\begin{lemma}
\label{hypercovering-lemma-h0-cech}
设 $\mathcal{C}$ 为有纤维积的位点，$X$ 为 $\mathcal{C}$ 的对象，
$K$ 为 $X$ 的超覆盖，且 $\mathcal{F}$ 为 $\mathcal{C}$ 上的阿贝尔群层。
则 $\check{H}^0(K, \mathcal{F}) = \mathcal{F}(X)$。
\end{lemma}

\begin{proof}
有
$$
\check{H}^0(K, \mathcal{F})
=
\Ker(\mathcal{F}(K_0) \longrightarrow \mathcal{F}(K_1))
$$
写 $K_0 = \{U_i \to X\}$。它是位点 $\mathcal{C}$ 中的覆盖。
此外，$K_1 \to K_0 \times K_0$ 是
$\text{SR}(\mathcal{C}, X)$ 中的覆盖。因此可写
$K_1 = \amalg_{i_0, i_1 \in I} \{V_{i_0i_1j} \to X\}$，
使态射 $K_1 \to K_0 \times K_0$ 由位点 $\mathcal{C}$ 的覆盖
$\{V_{i_0i_1j} \to U_{i_0} \times_X U_{i_1}\}$ 给出。
从而还可把
$$
\check{H}^0(K, \mathcal{F})
=
\Ker(
\prod\nolimits_i \mathcal{F}(U_i)
\longrightarrow
\prod\nolimits_{i_0i_1 j} \mathcal{F}(V_{i_0i_1j})
)
$$
按显然的映射作此等同。$\mathcal{F}$ 的层性质蕴含
$\check{H}^0(K, \mathcal{F}) = H^0(X, \mathcal{F})$。
\end{proof}

\noindent
事实上，该性质当然刻画了 $\mathcal{C}$ 上所有阿贝尔预层中的阿贝尔层。
在此情形下，《层上同调》引理 \ref{hypercovering-lemma-injective-trivial-cech}
的类似物如下。

\begin{lemma}
\label{hypercovering-lemma-injective-trivial-cech}
设 $\mathcal{C}$ 是具有纤维积的位点，$X$ 是 $\mathcal{C}$ 的对象，
$K$ 是 $X$ 的超覆盖，$\mathcal{I}$ 是 $\mathcal{C}$ 上的内射阿贝尔群层。则
$$
\check{H}^p(K, \mathcal{I}) =
\left\{
\begin{matrix}
\mathcal{I}(X) & \text{若} & p = 0 \\
0 & \text{若} & p > 0
\end{matrix}
\right.
$$
\end{lemma}

\begin{proof}
注意，对 $\text{SR}(\mathcal{C}, X)$ 的任意对象
$Z = \{U_i \to X\}$ 以及 $\mathcal{C}$ 上的任意阿贝尔层
$\mathcal{F}$，都有
\begin{eqnarray*}
\mathcal{F}(Z)
& = &
\prod \mathcal{F}(U_i) \\
& = &
\prod \Mor_{\textit{PSh}(\mathcal{C})}(h_{U_i}, \mathcal{F})\\
& = &
\Mor_{\textit{PSh}(\mathcal{C})}(F(Z), \mathcal{F})\\
& = &
\Mor_{\textit{PAb}(\mathcal{C})}(\mathbf{Z}_{F(Z)}, \mathcal{F}) \\
& = &
\Mor_{\textit{Ab}(\mathcal{C})}(\mathbf{Z}_{F(Z)}^\#, \mathcal{F})
\end{eqnarray*}
由此可见，对 $\text{SR}(\mathcal{C}, X)$ 的任意单纯对象 $K$，都有
\begin{equation}
\label{hypercovering-equation-identify-cech}
s(\mathcal{F}(K))
=
\Hom_{\textit{Ab}(\mathcal{C})}(s(\mathbf{Z}_{F(K)}^\#), \mathcal{F})
\end{equation}
记号见定义 \ref{hypercovering-definition-homology}。
若 $K$ 是超覆盖，则层复形 $s(\mathbf{Z}_{F(K)}^\#)$
拟同构于 $\mathbf{Z}_X^\#$，见引理 \ref{hypercovering-lemma-hypercovering-acyclic}。
因此，若 $\mathcal{I}$ 是内射阿贝尔层而 $K$ 是超覆盖，
则复形 $s(\mathcal{I}(K))$ 除了次数 $0$ 外是无环的。换言之，
$$
\check{H}^i(K, \mathcal{I}) = 0
$$
对 $i > 0$ 成立。结合引理 \ref{hypercovering-lemma-h0-cech} 即得结论。
\end{proof}

\noindent
下面给出《位点上的上同调》引理
\ref{sites-cohomology-lemma-cech-spectral-sequence} 的类似物。
设 $\mathcal{C}$ 是位点，$\mathcal{F}$ 是 $\mathcal{C}$ 上的阿贝尔群层。
回忆，$\underline{H}^i(\mathcal{F})$ 表示 $\mathcal{C}$ 上由规则
$\underline{H}^i(\mathcal{F}) : U \longmapsto H^i(U, \mathcal{F})$
定义的阿贝尔群预层。依照本节引言中的做法，把它扩张到
$\text{SR}(\mathcal{C})$ 上。

\begin{lemma}
\label{hypercovering-lemma-cech-spectral-sequence}
设 $\mathcal{C}$ 是具有纤维积的位点，$X$ 是 $\mathcal{C}$ 的对象，
$K$ 是 $X$ 的超覆盖，$\mathcal{F}$ 是 $\mathcal{C}$ 上的阿贝尔群层。
在 $D^{+}(\textit{Ab})$ 中存在关于 $\mathcal{F}$ 函子性的态射
$$
s(\mathcal{F}(K))
\longrightarrow
R\Gamma(X, \mathcal{F})
$$
它诱导函子 $\textit{Ab}(\mathcal{C}) \to \textit{Ab}$ 之间的自然变换
$$
\check{H}^i(K, -) \longrightarrow H^i(X, -)
$$
此外，存在谱序列 $(E_r, d_r)_{r \geq 0}$，满足
$$
E_2^{p, q} = \check{H}^p(K, \underline{H}^q(\mathcal{F}))
$$
并收敛到 $H^{p + q}(X, \mathcal{F})$。
此谱序列关于 $\mathcal{F}$ 和超覆盖 $K$ 都是函子性的。
\end{lemma}

\begin{proof}
可以采用《上同调》一章相应引理中的同一方法证明。这里改用双复形论证。

\medskip\noindent
在 $\mathcal{C}$ 上阿贝尔层的范畴中选取内射分辨
$\mathcal{F} \to \mathcal{I}^\bullet$。考虑各项为
$$
A^{p, q} = \mathcal{I}^q(K_p)
$$
的双复形 $A^{\bullet, \bullet}$，其中微分
$d_1^{p, q} : A^{p, q} \to A^{p + 1, q}$ 来自与余单纯阿贝尔群
$\mathcal{I}^p(K)$ 相伴的复形 $s(\mathcal{I}^q(K))$ 上的微分，
而微分 $d_2^{p, q} : A^{p, q} \to A^{p, q + 1}$ 来自微分
$\mathcal{I}^q \to \mathcal{I}^{q + 1}$。
用 $\text{Tot}(A^{\bullet, \bullet})$ 表示与双复形
$A^{\bullet, \bullet}$ 相伴的全复形，见《同调代数》第
\ref{homology-section-double-complexes} 节。我们将使用与此双复形相伴的
两个谱序列 $({}'E_r, {}'d_r)$ 和 $({}''E_r, {}''d_r)$，见
《同调代数》第 \ref{homology-section-double-complex} 节。

\medskip\noindent
由引理 \ref{hypercovering-lemma-injective-trivial-cech}，复形
$s(\mathcal{I}^q(K))$ 在正次数无环，且其 $H^0$ 等于
$\mathcal{I}^q(X)$。因此，由《同调代数》引理
\ref{homology-lemma-double-complex-gives-resolution}，自然映射
$$
\mathcal{I}^\bullet(X) \longrightarrow \text{Tot}(A^{\bullet, \bullet})
$$
是阿贝尔群复形的拟同构。特别地，
$H^n(\text{Tot}(A^{\bullet, \bullet})) = H^n(X, \mathcal{F})$。

\medskip\noindent
引理中的态射 $s(\mathcal{F}(K)) \longrightarrow R\Gamma(X, \mathcal{F})$
是态射 $s(\mathcal{F}(K)) \to \text{Tot}(A^{\bullet, \bullet})$
与上述拟同构之逆的复合。这是可行的，因为 $\mathcal{I}^\bullet(X)$
是 $R\Gamma(X, \mathcal{F})$ 的一个代表元。

\medskip\noindent
考虑谱序列 $({}'E_r, {}'d_r)_{r \geq 0}$。由《同调代数》引理
\ref{homology-lemma-ss-double-complex}，有
$$
{}'E_2^{p, q} = H^p_I(H^q_{II}(A^{\bullet, \bullet}))
$$
换言之，先对 $d_2$ 取上同调，得到群
${}'E_1^{p, q} = \underline{H}^q(\mathcal{F})(K_p)$。
因此，由微分 ${}'d_1$ 的描述，确有
${}'E_2^{p, q} = \check{H}^p(K, \underline{H}^q(\mathcal{F}))$。
由以上所述和《同调代数》引理
\ref{homology-lemma-first-quadrant-ss}，可见它按要求收敛到
$H^n(X, \mathcal{F})$。

\medskip\noindent
关于上述构造对阿贝尔层 $\mathcal{F}$ 和超覆盖 $K$ 的函子性，
其证明从略。
\end{proof}









\section{Verdier 意义下的超覆盖}
\label{hypercovering-section-hypercoverings-verdier}

\noindent
敏锐的读者会注意到，为了得到对象 $X$ 的超覆盖所对应的从
{\v C}ech 上同调到上同调的谱序列，我们所需的只是引理
\ref{hypercovering-lemma-hypercovering-F} 的结论。因此下面的定义是合理的。

\begin{definition}
\label{hypercovering-definition-hypercovering-variant}
设 $\mathcal{C}$ 是位点，并假设 $\mathcal{C}$ 具有等化子和纤维积。
设 $\mathcal{G}$ 是集合预层。$\mathcal{G}$ 的一个{\it 超覆盖}，是
$\text{SR}(\mathcal{C})$ 的单纯对象 $K$ 连同一个增广
$F(K) \to \mathcal{G}$，使得
\begin{enumerate}
\item $F(K_0) \to \mathcal{G}$ 在层化后成为满射；
\item $F(K_1) \to F(K_0) \times_\mathcal{G} F(K_0)$
在层化后成为满射；并且
\item $F(K_{n + 1}) \longrightarrow F((\text{cosk}_n \text{sk}_n K)_{n + 1})$
对 $n \geq 1$ 在层化后成为满射。
\end{enumerate}
若 $\text{SR}(\mathcal{C})$ 的单纯对象 $K$ 是
$\textit{PSh}(\mathcal{C})$ 的终对象 $*$ 的超覆盖，则称 $K$ 为一个
{\it 超覆盖}。
\end{definition}

\noindent
假设 $\mathcal{C}$ 具有纤维积和等化子，便保证
$\text{SR}(\mathcal{C})$ 具有纤维积和等化子，并且 $F$ 与它们交换
（引理 \ref{hypercovering-lemma-coprod-prod-SR}）；这足以定义所用的余骨架函子
（见《单纯方法》评注 \ref{simplicial-remark-existence-cosk} 和
《范畴论》引理 \ref{categories-lemma-fibre-products-equalizers-exist}）。
对一般的 $\mathcal{C}$，可把条件 (3) 替换为：
$F(K_{n + 1}) \longrightarrow ((\text{cosk}_n \text{sk}_n F(K))_{n + 1})$
对 $n \geq 1$ 在层化后成为满射；本节的结果仍然成立。

\medskip\noindent
设 $\mathcal{F}$ 是 $\mathcal{C}$ 上的阿贝尔层。上一节定义了
$\mathcal{F}$ 关于 $\text{SR}(\mathcal{C})$ 的单纯对象 $K$ 的
{\v C}ech 复形。现在，对给定预层 $\mathcal{G}$，置
$$
H^0(\mathcal{G}, \mathcal{F}) =
\Mor_{\textit{PSh}(\mathcal{C})}(\mathcal{G}, \mathcal{F}) =
\Mor_{\Sh(\mathcal{C})}(\mathcal{G}^\#, \mathcal{F}) =
H^0(\mathcal{G}^\#, \mathcal{F})
$$
记号同《位点上的上同调》第 \ref{sites-cohomology-section-limp} 节。
这是左正合函子，其高阶导出函子（《位点上的上同调》第
\ref{sites-cohomology-section-limp} 节作了简要研究）记为
$H^i(\mathcal{G}, \mathcal{F})$。我们将证明，给定 $\mathcal{G}$ 的
超覆盖 $K$，存在一个从 {\v C}ech 上同调到上同调的谱序列，
收敛到上同调 $H^i(\mathcal{G}, \mathcal{F})$。注意，若
$\mathcal{G} = *$，则 $H^i(*, \mathcal{F}) = H^i(\mathcal{C}, \mathcal{F})$，
这恢复了 $\mathcal{F}$ 在位点 $\mathcal{C}$ 上的上同调。

\begin{lemma}
\label{hypercovering-lemma-h0-cech-variant}
设 $\mathcal{C}$ 是具有等化子和纤维积的位点，
$\mathcal{G}$ 是 $\mathcal{C}$ 上的预层，$K$ 是 $\mathcal{G}$ 的超覆盖，
$\mathcal{F}$ 是 $\mathcal{C}$ 上的阿贝尔群层。则
$\check{H}^0(K, \mathcal{F}) = H^0(\mathcal{G}, \mathcal{F})$。
\end{lemma}

\begin{proof}
这由 $H^0(\mathcal{G}, \mathcal{F})$ 的定义以及下图在层化后成为余等化子图的事实得出：
$$
\xymatrix{
F(K_1) \ar@<1ex>[r] \ar@<-1ex>[r] &
F(K_0) \ar[r] & \mathcal{G}
}
$$
\end{proof}

\begin{lemma}
\label{hypercovering-lemma-injective-trivial-cech-variant}
设 $\mathcal{C}$ 是具有等化子和纤维积的位点，
$\mathcal{G}$ 是 $\mathcal{C}$ 上的预层，$K$ 是 $\mathcal{G}$ 的超覆盖，
$\mathcal{I}$ 是 $\mathcal{C}$ 上的内射阿贝尔群层。则
$$
\check{H}^p(K, \mathcal{I}) =
\left\{
\begin{matrix}
H^0(\mathcal{G}, \mathcal{I}) & \text{若} & p = 0 \\
0 & \text{若} & p > 0
\end{matrix}
\right.
$$
\end{lemma}

\begin{proof}
由 (\ref{hypercovering-equation-identify-cech})，有
$$
s(\mathcal{F}(K))
=
\Hom_{\textit{Ab}(\mathcal{C})}(s(\mathbf{Z}_{F(K)}^\#), \mathcal{F})
$$
复形 $s(\mathbf{Z}_{F(K)}^\#)$ 拟同构于
$\mathbf{Z}_\mathcal{G}^\#$，见引理
\ref{hypercovering-lemma-acyclic-hypercover-sheaves}。因此，若 $\mathcal{I}$ 是内射阿贝尔层，
则复形 $s(\mathcal{I}(K))$ 除了次数 $0$ 外是无环的。换言之，
对 $i > 0$ 有 $\check{H}^i(K, \mathcal{I}) = 0$。
结合引理 \ref{hypercovering-lemma-h0-cech-variant} 即得结论。
\end{proof}

\begin{lemma}
\label{hypercovering-lemma-cech-spectral-sequence-variant}
设 $\mathcal{C}$ 是具有等化子和纤维积的位点，
$\mathcal{G}$ 是 $\mathcal{C}$ 上的预层，$K$ 是 $\mathcal{G}$ 的超覆盖，
$\mathcal{F}$ 是 $\mathcal{C}$ 上的阿贝尔群层。
在 $D^{+}(\textit{Ab})$ 中存在关于 $\mathcal{F}$ 函子性的态射
$$
s(\mathcal{F}(K)) \longrightarrow R\Gamma(\mathcal{G}, \mathcal{F})
$$
它诱导函子 $\textit{Ab}(\mathcal{C}) \to \textit{Ab}$ 之间的自然变换
$$
\check{H}^i(K, -) \longrightarrow H^i(\mathcal{G}, -)
$$
此外，存在谱序列 $(E_r, d_r)_{r \geq 0}$，满足
$$
E_2^{p, q} = \check{H}^p(K, \underline{H}^q(\mathcal{F}))
$$
并收敛到 $H^{p + q}(\mathcal{G}, \mathcal{F})$。
此谱序列关于 $\mathcal{F}$ 和超覆盖 $K$ 都是函子性的。
\end{lemma}

\begin{proof}
在 $\mathcal{C}$ 上阿贝尔层的范畴中选取内射分辨
$\mathcal{F} \to \mathcal{I}^\bullet$。考虑各项为
$$
A^{p, q} = \mathcal{I}^q(K_p)
$$
的双复形 $A^{\bullet, \bullet}$，其中微分
$d_1^{p, q} : A^{p, q} \to A^{p + 1, q}$ 来自微分
$\mathcal{I}^p \to \mathcal{I}^{p + 1}$，而微分
$d_2^{p, q} : A^{p, q} \to A^{p, q + 1}$ 来自与余单纯阿贝尔群
$\mathcal{I}^p(K)$ 相伴的复形 $s(\mathcal{I}^p(K))$ 上的微分，
如上所述。我们将使用与此双复形相伴的两个谱序列
$({}'E_r, {}'d_r)$ 和 $({}''E_r, {}''d_r)$，见《同调代数》第
\ref{homology-section-double-complex} 节。

\medskip\noindent
由引理 \ref{hypercovering-lemma-injective-trivial-cech-variant}，复形
$s(\mathcal{I}^p(K))$ 在正次数无环，且其 $H^0$ 等于
$H^0(\mathcal{G}, \mathcal{I}^p)$。因此，由《同调代数》引理
\ref{homology-lemma-double-complex-gives-resolution} 及其证明，
谱序列 $({}'E_r, {}'d_r)$ 退化，并且自然映射
$$
H^0(\mathcal{G}, \mathcal{I}^\bullet) \longrightarrow
\text{Tot}(A^{\bullet, \bullet})
$$
是阿贝尔群复形的拟同构。引理中的态射
$s(\mathcal{F}(K)) \longrightarrow R\Gamma(\mathcal{G}, \mathcal{F})$
是自然映射 $s(\mathcal{F}(K)) \to \text{Tot}(A^{\bullet, \bullet})$
与上述拟同构之逆的复合。这是可行的，因为
$H^0(\mathcal{G}, \mathcal{I}^\bullet)$ 是
$R\Gamma(\mathcal{G}, \mathcal{F})$ 的一个代表元。

\medskip\noindent
考虑谱序列 $({}''E_r, {}''d_r)_{r \geq 0}$。由《同调代数》引理
\ref{homology-lemma-ss-double-complex}，有
$$
{}''E_2^{p, q} = H^p_{II}(H^q_I(A^{\bullet, \bullet}))
$$
换言之，先对 $d_1$ 取上同调，得到群
${}''E_1^{p, q} = \underline{H}^p(\mathcal{F})(K_q)$。
因此，由微分 ${}''d_1$ 的描述，确有
${}''E_2^{p, q} = \check{H}^p(K, \underline{H}^q(\mathcal{F}))$。
由于此谱序列收敛到 $\text{Tot}(A^{\bullet, \bullet})$ 的上同调，
证明完成。
\end{proof}

\begin{lemma}
\label{hypercovering-lemma-cech-spectral-sequence-verdier}
设 $\mathcal{C}$ 是具有等化子和纤维积的位点，$K$ 是超覆盖，
$\mathcal{F}$ 是阿贝尔层。存在谱序列 $(E_r, d_r)_{r \geq 0}$，满足
$$
E_2^{p, q} = \check{H}^p(K, \underline{H}^q(\mathcal{F}))
$$
并收敛到整体上同调群 $H^{p + q}(\mathcal{F})$。
\end{lemma}

\begin{proof}
这是引理 \ref{hypercovering-lemma-cech-spectral-sequence-variant} 的特殊情形。
\end{proof}







\section{超覆盖的覆盖细化}
\label{hypercovering-section-covering}

\noindent
下面给出若干构造超覆盖的方法。注意，由于范畴
$\text{SR}(\mathcal{C}, X)$ 具有纤维积，它的单纯对象范畴也具有纤维积，
见《单纯方法》引理 \ref{simplicial-lemma-fibre-product}。

\begin{lemma}
\label{hypercovering-lemma-funny-gamma}
设 $\mathcal{C}$ 是具有纤维积的位点，$X$ 是 $\mathcal{C}$ 的对象，
$K,L,M$ 是 $\text{SR}(\mathcal{C}, X)$ 的单纯对象，并设
$a : K \to L$、$b : M \to L$ 是态射。假设
\begin{enumerate}
\item $K$ 是 $X$ 的超覆盖；
\item 态射 $M_0 \to L_0$ 是覆盖；并且
\item 对每个 $n \geq 0$，下图中的箭头 $\gamma$ 是覆盖：
$$
\xymatrix{
M_{n + 1} \ar[dd] \ar[rr] \ar[rd]^\gamma &
&
(\text{cosk}_n \text{sk}_n M)_{n + 1} \ar[dd] \\
&
L_{n + 1}
\times_{(\text{cosk}_n \text{sk}_n L)_{n + 1}}
(\text{cosk}_n \text{sk}_n M)_{n + 1}
\ar[ld] \ar[ru]
& \\
L_{n + 1} \ar[rr] & & (\text{cosk}_n \text{sk}_n L)_{n + 1}
}
$$
\end{enumerate}
则纤维积 $K \times_L M$ 是 $X$ 的超覆盖。
\end{lemma}

\begin{proof}
由 (2)，态射 $(K \times_L M)_0 = K_0 \times_{L_0} M_0 \to K_0$
是覆盖的基变换，因而是覆盖，见引理 \ref{hypercovering-lemma-covering-permanence}。
由 (1)，$K_0 \to \{X \to X\}$ 是覆盖。因此，由引理
\ref{hypercovering-lemma-covering-permanence}，
$(K \times_L M)_0 \to \{X \to X\}$ 是覆盖。故
$K \times_L M$ 满足定义 \ref{hypercovering-definition-hypercovering} 的第一个条件。

\medskip\noindent
还需验证，对所有 $n \geq 0$，态射
$$
K_{n + 1} \times_{L_{n + 1}} M_{n + 1} = (K \times_L M)_{n + 1}
\longrightarrow
(\text{cosk}_n \text{sk}_n (K \times_L M))_{n + 1}
$$
是覆盖。简记
$A = (\text{cosk}_n \text{sk}_n K)_{n + 1}$,
$B = (\text{cosk}_n \text{sk}_n L)_{n + 1}$，以及
$C = (\text{cosk}_n \text{sk}_n M)_{n + 1}$.
函子 $\text{cosk}_n \text{sk}_n$ 与纤维积交换，见《单纯方法》引理
\ref{simplicial-lemma-cosk-fibre-product}。因此，上式右端等于
$A \times_B C$。考虑交换图
$$
\xymatrix{
K_{n + 1} \times_{L_{n + 1}} M_{n + 1} \ar[r] \ar[d] &
M_{n + 1} \ar[d] \ar[rd]_\gamma \ar[rrd] &
& \\
K_{n + 1} \ar[r] \ar[rd] &
L_{n + 1} \ar[rrd] &
L_{n + 1} \times_B C \ar[l] \ar[r] &
C \ar[d] \\
&
A \ar[rr] &
&
B
}
$$
此图表明
$$
K_{n + 1} \times_{L_{n + 1}} M_{n + 1}
=
(K_{n + 1} \times_B C)
\times_{(L_{n + 1} \times_B C), \gamma}
M_{n + 1}
$$
现在，$K_{n + 1} \times_B C \to A \times_B C$ 是覆盖
$K_{n + 1} \to A$ 沿态射 $A \times_B C \to A$ 的基变换，
故为覆盖。由假设 (3)，态射 $\gamma$ 是覆盖。因此态射
$$
(K_{n + 1} \times_B C)
\times_{(L_{n + 1} \times_B C), \gamma}
M_{n + 1}
\longrightarrow
K_{n + 1} \times_B C
$$
作为覆盖的基变换也是覆盖。覆盖的复合仍是覆盖，故引理成立。
\end{proof}

\begin{lemma}
\label{hypercovering-lemma-product-hypercoverings}
设 $\mathcal{C}$ 是具有纤维积的位点，$X$ 是 $\mathcal{C}$ 的对象。
若 $K,L$ 是 $X$ 的超覆盖，则 $K \times L$ 是 $X$ 的超覆盖。
\end{lemma}

\begin{proof}
既可直接验证，也可应用上述引理 \ref{hypercovering-lemma-funny-gamma}，并验证
$L \to \{X \to X\}$ 具有性质 (3)。
\end{proof}


\noindent
设 $\mathcal{C}$ 是具有纤维积的位点，$X$ 是 $\mathcal{C}$ 的对象。
由于范畴 $\text{SR}(\mathcal{C}, X)$ 具有余积和有限极限，
对某些单纯集 $U$（例如只有有限多个非退化单形者）以及
$\text{SR}(\mathcal{C}, X)$ 的任意单纯对象 $K$，可以讨论对象
$U \times K$ 和 $\Hom(U, K)$。见《单纯方法》第
\ref{simplicial-section-product-with-simplicial-sets} 节和第
\ref{simplicial-section-hom-from-simplicial-sets} 节。

\begin{lemma}
\label{hypercovering-lemma-covering}
设 $\mathcal{C}$ 是具有纤维积的位点，$X$ 是 $\mathcal{C}$ 的对象，
$K$ 是 $X$ 的超覆盖，$k \geq 0$ 是整数。设
$u : Z \to K_k$ 是 $\text{SR}(\mathcal{C}, X)$ 中的覆盖。
则存在超覆盖的态射 $f: L \to K$，使得 $L_k \to K_k$ 经由 $u$ 分解。
\end{lemma}

\begin{proof}
记 $Y = K_k$。设 $C[k]$ 是《单纯方法》例
\ref{simplicial-example-simplex-cosimplicial-set} 中定义的余单纯集。
我们将使用《单纯方法》引理
\ref{simplicial-lemma-morphism-into-product} 给出的
$\Hom(C[k], Y)$ 和 $\Hom(C[k], Z)$ 的描述。
存在对应于 $\text{id} : K_k = Y \to Y$ 的典范态射
$K \to \Hom(C[k], Y)$。考虑态射
$\Hom(C[k], Z) \to \Hom(C[k], Y)$；其次数 $n$ 项上的态射为
$$
\prod\nolimits_{\alpha : [k] \to [n]} Z
\longrightarrow
\prod\nolimits_{\alpha : [k] \to [n]} Y
$$
即在每个因子上使用给定态射 $Z \to Y$。置
$$
L = K \times_{\Hom(C[k], Y)} \Hom(C[k], Z).
$$
态射 $L_k \to K_k$ 位于交换图
$$
\xymatrix{
L_k \ar[r] \ar[d] &
\prod_{\alpha : [k] \to [k]} Z \ar[r]^-{\text{pr}_{\text{id}_{[k]}}} \ar[d] &
Z \ar[d] \\
K_k \ar[r] &
\prod_{\alpha : [k] \to [k]} Y \ar[r]^-{\text{pr}_{\text{id}_{[k]}}} &
Y
}
$$
中。由于底部两个箭头的复合是恒等态射，便得到所需分解。

\medskip\noindent
还需证明 $L$ 是 $X$ 的超覆盖。为此应用引理
\ref{hypercovering-lemma-funny-gamma}。条件 (1) 由假设成立。对于 (2)，态射
$$
\Hom(C[k], Z)_0 \to \Hom(C[k], Y)_0
$$
是覆盖，因为态射 $[k] \to [0]$ 仅有一个，故它同构于 $Z \to Y$。

\medskip\noindent
考虑 $n = 0$ 时的条件 (3)。由于
$(\text{cosk}_0 T)_1 = T \times T$
（《单纯方法》例 \ref{simplicial-example-cosk0}）且
$\Hom(C[k], Z)_1 = \prod_{\alpha : [k] \to [1]} Z$，得到图
$$
\xymatrix{
\prod\nolimits_{\alpha : [k] \to [1]} Z \ar[r] \ar[d] &
Z \times Z \ar[d] \\
\prod\nolimits_{\alpha : [k] \to [1]} Y \ar[r] &
Y \times Y
}
$$
其中水平箭头对应于向两个非满射 $\alpha$ 所对应因子的投影。
因此箭头 $\gamma$ 是态射
$$
\prod\nolimits_{\alpha : [k] \to [1]} Z
\longrightarrow
\prod\nolimits_{\alpha : [k] \to [1]\text{ 非满}} Z
\times
\prod\nolimits_{\alpha : [k] \to [1]\text{ 满}} Y
$$
它是覆盖的乘积，故由引理 \ref{hypercovering-lemma-covering-permanence} 仍为覆盖。

\medskip\noindent
考虑 $n > 0$ 时的条件 (3)。我们断言，存在有限集的单射
$\tau : S' \to S$，使得对 $\text{SR}(\mathcal{C}, X)$ 的任意对象
$T$，态射
\begin{equation}
\label{hypercovering-equation-map}
\Hom(C[k], T)_{n + 1} \to
(\text{cosk}_n \text{sk}_n \Hom(C[k], T))_{n + 1}
\end{equation}
关于 $T$ 函子性地同构于投影
$\prod_{s \in S} T \to \prod_{s' \in S'} T$。若此言成立，
则仿照上一段可见，箭头 $\gamma$ 是态射
$$
\prod\nolimits_{s \in S} Z
\longrightarrow
\prod\nolimits_{s \in S'} Z
\times
\prod\nolimits_{s \not\in \tau(S')} Y
$$
它是覆盖的乘积，故由引理 \ref{hypercovering-lemma-covering-permanence} 仍为覆盖。
由构造，有
$\Hom(C[k], T)_{n + 1} = \prod_{\alpha : [k] \to [n + 1]} T$
（见《单纯方法》引理 \ref{simplicial-lemma-morphism-into-product}）。
相应地取 $S = \text{Map}([k], [n + 1])$。另一方面，
《单纯方法》引理 \ref{simplicial-lemma-formula-limit} 把
$(\text{cosk}_n \text{sk}_n \Hom(C[k], T))_{n + 1}$
的点描述为 $\Hom(C[k], T)_n$ 的点列
$(f_0, \ldots, f_{n + 1})$，满足
$d^n_{j - 1} f_i = d^n_i f_j$（$0 \leq i < j \leq n + 1$）。
可写 $f_i = (f_{i, \alpha})$，其中 $f_{i, \alpha}$ 是 $T$ 的点，
$\alpha \in \text{Map}([k], [n])$。这些条件化为
$$
f_{i, \delta^n_{j - 1} \circ \beta} = f_{j, \delta_i^n \circ \beta}
$$
对任意 $0 \leq i < j \leq n + 1$ 和 $\beta : [k] \to [n - 1]$
成立。因此
$$
S' = \{0, \ldots, n + 1\} \times \text{Map}([k], [n]) / \sim
$$
其中等价关系由等价式
$$
(i, \delta^n_{j - 1} \circ \beta) \sim (j, \delta_i^n \circ \beta)
$$
生成，其中 $0 \leq i < j \leq n + 1$ 且
$\beta : [k] \to [n - 1]$。一个略去的计算表明，态射
(\ref{hypercovering-equation-map}) 对应于映射 $S' \to S$，它把 $(i, \alpha)$
送到 $\delta^{n + 1}_i \circ \alpha \in S$。（由《单纯方法》引理
\ref{simplicial-lemma-relations-face-degeneracy} 的 (1)，此映射确实良定义，
这也许能使读者安心。）为完成证明，只需说明：若
$\alpha, \alpha' : [k] \to [n]$、$0 \leq i < j \leq n + 1$，且
$$
\delta^{n + 1}_i \circ \alpha = \delta^{n + 1}_j \circ \alpha'
$$
则存在某个 $\beta : [k] \to [n - 1]$，使得
$\alpha = \delta^n_{j - 1} \circ \beta$ 且
$\alpha' = \delta_i^n \circ \beta$。这很容易验证，故略去。
\end{proof}

\begin{lemma}
\label{hypercovering-lemma-covering-sheaf}
设 $\mathcal{C}$ 是具有纤维积的位点，$X$ 是 $\mathcal{C}$ 的对象，
$K$ 是 $X$ 的超覆盖，$n \geq 0$ 是整数。设
$u : \mathcal{F} \to F(K_n)$ 是层化后成为满射的预层态射。
则存在超覆盖的态射 $f: L \to K$，使得
$F(f_n) : F(L_n) \to F(K_n)$ 经由 $u$ 分解。
\end{lemma}

\begin{proof}
写成 $K_n = \{U_i \to X\}_{i \in I}$。于是映射 $u$ 是集合预层态射
$u : \mathcal{F} \to \amalg h_{u_i}$。关于 $u$ 的假设意味着：
对每个 $i \in I$，存在位点 $\mathcal{C}$ 的覆盖
$\{U_{ij} \to U_i\}_{j \in I_i}$ 和预层态射
$t_{ij} : h_{U_{ij}} \to \mathcal{F}$，使得 $u \circ t_{ij}$
是由态射 $U_{ij} \to U_i$ 得到的映射
$h_{U_{ij}} \to h_{U_i}$。置 $J = \amalg_{i \in I} I_i$，并令
$\alpha : J \to I$ 为显然的映射。对 $j \in J$，记
$V_j = U_{\alpha(j)j}$，并置 $Z = \{V_j \to X\}_{j \in J}$。
最后，考虑由 $\alpha : J \to I$ 和上述态射
$V_j = U_{\alpha(j)j} \to U_{\alpha(j)}$ 给出的态射
$u' : Z \to K_n$。显然，这是范畴 $\text{SR}(\mathcal{C}, X)$ 中的覆盖；
由构造，$F(u') : F(Z) \to F(K_n)$ 经由 $u$ 分解。
故结论由上述引理 \ref{hypercovering-lemma-covering} 得出。
\end{proof}


\section{添加单形}
\label{hypercovering-section-adding-simplices}

\noindent
本节证明后面将用到的一些技术性引理。设 $\mathcal{C}$ 是具有纤维积的位点，
$X$ 是 $\mathcal{C}$ 的对象。如上述第 \ref{hypercovering-section-covering} 节所指出，
对某些单纯集 $U$ 和 $\text{SR}(\mathcal{C}, X)$ 的任意单纯对象 $K$，
对象 $U \times K$ 和 $\Hom(U, K)$ 都已定义。见《单纯方法》第
\ref{simplicial-section-product-with-simplicial-sets} 节和第
\ref{simplicial-section-hom-from-simplicial-sets} 节。

\begin{lemma}
\label{hypercovering-lemma-one-more-simplex}
设 $\mathcal{C}$ 是具有纤维积的位点，$X$ 是 $\mathcal{C}$ 的对象，
$K$ 是 $X$ 的超覆盖。设 $U \subset V$ 是单纯集，且对所有 $n$，
$U_n,V_n$ 都是有限非空集。假设 $U$ 只有有限多个非退化单形。
设 $n \geq 0$ 且 $x \in V_n$、$x \not \in U_n$，并满足
\begin{enumerate}
\item 对 $i < n$ 有 $V_i = U_i$；
\item $V_n = U_n \cup \{x\}$；
\item 对 $j > n$，任意 $z \in V_j$、$z \not \in U_j$ 都是退化的。
\end{enumerate}
则 $\text{SR}(\mathcal{C}, X)$ 中的态射
$$
\Hom(V, K)_0
\longrightarrow
\Hom(U, K)_0
$$
是覆盖。
\end{lemma}

\begin{proof}
若 $n = 0$，则容易推出 $V = U \amalg \Delta[0]$（见下文）。
此时 $\Hom(V, K)_0 = \Hom(U, K)_0 \times K_0$，结论由引理
\ref{hypercovering-lemma-covering-permanence} 得出。

\medskip\noindent
设 $a : \Delta[n] \to V$ 是《单纯方法》引理
\ref{simplicial-lemma-simplex-map} 中与 $x$ 相伴的态射。
记 $\partial \Delta[n] = i_{(n-1)!} \text{sk}_{n - 1} \Delta[n]$
为 $\Delta[n]$ 的 $(n - 1)$-骨架。设
$b : \partial \Delta[n] \to U$ 是 $a$ 在 $\Delta[n]$ 的
$(n - 1)$-骨架上的限制。由《单纯方法》引理
\ref{simplicial-lemma-glue-simplex}，有
$V = U \amalg_{\partial \Delta[n]} \Delta[n]$。由《单纯方法》引理
\ref{simplicial-lemma-hom-from-coprod}，图
$$
\xymatrix{
\Hom(V, K)_0 \ar[r] \ar[d] &
\Hom(U, K)_0 \ar[d] \\
\Hom(\Delta[n], K)_0 \ar[r] &
\Hom(\partial \Delta[n], K)_0
}
$$
是纤维积方块。因此只需证明底部水平箭头是覆盖。由《单纯方法》引理
\ref{simplicial-lemma-cosk-shriek}，此箭头等同于
$$
K_n \to (\text{cosk}_{n - 1} \text{sk}_{n - 1} K)_n
$$
故由超覆盖的定义，它是覆盖。
\end{proof}

\begin{lemma}
\label{hypercovering-lemma-add-simplices}
设 $\mathcal{C}$ 是具有纤维积的位点，$X$ 是 $\mathcal{C}$ 的对象，
$K$ 是 $X$ 的超覆盖。设 $U \subset V$ 是单纯集，且对所有 $n$，
$U_n,V_n$ 都是有限非空集。假设 $U$ 和 $V$ 都只有有限多个非退化单形。
则 $\text{SR}(\mathcal{C}, X)$ 中的态射
$$
\Hom(V, K)_0
\longrightarrow
\Hom(U, K)_0
$$
是覆盖。
\end{lemma}

\begin{proof}
由上述引理 \ref{hypercovering-lemma-one-more-simplex}，只需证明关于引理中这类单纯集
包含 $U \subset V$ 的一个简单引理。这恰是《单纯方法》引理
\ref{simplicial-lemma-add-simplices} 的结论。
\end{proof}

\begin{lemma}
\label{hypercovering-lemma-degeneracy-maps-coverings}
设 $\mathcal{C}$ 是具有纤维积的位点，$X$ 是 $\mathcal{C}$ 的对象，
$K$ 是 $X$ 的超覆盖。则
\begin{enumerate}
\item 对每个 $n \geq 0$，$K_n$ 是 $X$ 的覆盖；
\item 对所有 $n \geq 1$ 和 $0 \leq i \leq n$，
$d^n_i : K_n \to K_{n - 1}$ 是覆盖。
\end{enumerate}
\end{lemma}

\begin{proof}
回忆，由定义 \ref{hypercovering-definition-hypercovering}，$K_0$ 是 $X$ 的覆盖；
这等价于说 $K_0 \to \{X \to X\}$ 是定义
\ref{hypercovering-definition-covering-SR} 意义下的覆盖。因此 (1) 由 (2) 得出，
因为 (2) 将证明复合
$K_n \to K_{n - 1} \to \ldots \to K_0 \to \{X \to X\}$
由引理 \ref{hypercovering-lemma-covering-permanence} 是覆盖。

\medskip\noindent
证明 (2)。由《单纯方法》引理
\ref{simplicial-lemma-exists-hom-from-simplicial-set-finite}，有
$\Mor(\Delta[n], K)_0 = K_n$。因此，把引理
\ref{hypercovering-lemma-add-simplices} 应用于 $n + 1$ 个不同的包含
$\Delta[n - 1] \to \Delta[n]$，便得 (2)。
\end{proof}

\begin{remark}
\label{hypercovering-remark-P-covering}
引理 \ref{hypercovering-lemma-add-simplices} 和
\ref{hypercovering-lemma-degeneracy-maps-coverings} 有如下有用的特殊情形。
设范畴 $\mathcal{C}$ 具有纤维积。设
$P \subset \text{Arrows}(\mathcal{C})$ 在基变换和复合下稳定，
并包含所有同构。所谓一个 {\it $P$-超覆盖}，是从
$\mathcal{C}$ 的单纯对象出发的增广 $a : U \to X$，使得
\begin{enumerate}
\item $U_0 \to X$ 属于 $P$；
\item $U_1 \to U_0 \times_X U_0$ 属于 $P$；
\item $U_{n + 1} \to (\text{cosk}_n\text{sk}_n U)_{n + 1}$
对 $n \geq 1$ 属于 $P$。
\end{enumerate}
范畴 $\mathcal{C}/X$ 具有所有有限极限，故上述表述中所用的余骨架存在
（见《范畴论》引理 \ref{categories-lemma-finite-limits-exist}）。
我们断言，态射 $U_n \to X$ 和 $d^n_i : U_n \to U_{n - 1}$ 都属于
$P$。把 $\mathcal{C}$ 视为一个位点，其覆盖为满足 $f \in P$ 的
$\{f : V \to U\}$，再取由 $K_n = \{U_n \to X\}$ 给出的 $K$，
便可由上述引理得到此结论。
\end{remark}


\section{同伦}
\label{hypercovering-section-homotopies}

\noindent
设 $\mathcal{C}$ 是具有纤维积的位点，$X$ 是 $\mathcal{C}$ 的对象，
$L$ 是 $\text{SR}(\mathcal{C}, X)$ 的单纯对象。由《单纯方法》引理
\ref{simplicial-lemma-exists-hom-from-simplicial-set-finite}，
范畴 $\text{Simp}(\text{SR}(\mathcal{C}, X))$ 中存在对象
$\Hom(\Delta[1], L)$，它表示函子
$$
T
\longmapsto
\Mor_{\text{Simp}(\text{SR}(\mathcal{C}, X))}(\Delta[1] \times T, L)
$$
由 $e_i : \Delta[0] \to \Delta[1]$ 和等同
$\Hom(\Delta[0], L) = L$，得到典范态射
$$
\Hom(\Delta[1], L) \to L \times L
$$

\begin{lemma}
\label{hypercovering-lemma-hom-hypercovering}
设 $\mathcal{C}$ 是具有纤维积的位点，$X$ 是 $\mathcal{C}$ 的对象，
$L$ 是 $\text{SR}(\mathcal{C}, X)$ 的单纯对象。设 $n \geq 0$。
考虑由上述态射得到的交换图
\begin{equation}
\label{hypercovering-equation-diagram}
\xymatrix{
\Hom(\Delta[1], L)_{n + 1} \ar[r] \ar[d] &
(\text{cosk}_n \text{sk}_n \Hom(\Delta[1], L))_{n + 1} \ar[d] \\
(L \times L)_{n + 1} \ar[r] &
(\text{cosk}_n \text{sk}_n (L \times L))_{n + 1}
}
\end{equation}
此图各项可作如下等同，其中
$\partial \Delta[n + 1] = i_{n!}\text{sk}_n \Delta[n + 1]$
是 $(n + 1)$-单形的 $n$-骨架：
\begin{eqnarray*}
\Hom(\Delta[1], L)_{n + 1}
& = &
\Hom(\Delta[1] \times \Delta[n + 1], L)_0 \\
(\text{cosk}_n \text{sk}_n \Hom(\Delta[1], L))_{n + 1}
& = &
\Hom(\Delta[1] \times \partial \Delta[n + 1], L)_0 \\
(L \times L)_{n + 1}
& = &
\Hom(
(\Delta[n + 1] \amalg \Delta[n + 1], L)_0 \\
(\text{cosk}_n \text{sk}_n (L \times L))_{n + 1}
& = &
\Hom(
\partial \Delta[n + 1]
\amalg
\partial \Delta[n + 1], L)_0
\end{eqnarray*}
而 $\text{SR}(\mathcal{C}, X)$ 中这些对象之间的态射来自单纯集交换图
\begin{equation}
\label{hypercovering-equation-dual-diagram}
\xymatrix{
\Delta[1] \times \Delta[n + 1] &
\Delta[1] \times \partial\Delta[n + 1] \ar[l] \\
\Delta[n + 1] \amalg \Delta[n + 1] \ar[u] &
\partial\Delta[n + 1] \amalg \partial\Delta[n + 1]
\ar[l] \ar[u]
}
\end{equation}
此外，图 (\ref{hypercovering-equation-diagram}) 中底部箭头与右侧箭头的纤维积等于
$$
\Hom(U, L)_0
$$
其中 $U \subset \Delta[1] \times \Delta[n + 1]$ 是最小单纯子集，
使得 $\Delta[n + 1] \amalg \Delta[n + 1]$ 和
$\Delta[1] \times \partial\Delta[n + 1]$ 都映入其中。
\end{lemma}

\begin{proof}
第一和第三个等式是《单纯方法》引理
\ref{simplicial-lemma-exists-hom-from-simplicial-set-finite}。
第二和第四个等式由该引理与《单纯方法》引理
\ref{simplicial-lemma-cosk-shriek} 合并得出。最后一个断言由如下事实得出：
$U$ 是图 (\ref{hypercovering-equation-dual-diagram}) 的底部和右侧箭头的推出；
再应用《单纯方法》引理 \ref{simplicial-lemma-hom-from-coprod}。
要看出 $U$ 等于此推出，只需看出
$\Delta[n + 1] \amalg \Delta[n + 1]$ 与
$\Delta[1] \times \partial\Delta[n + 1]$ 在
$\Delta[1] \times \Delta[n + 1]$ 中的交等于
$\partial\Delta[n + 1] \amalg \partial\Delta[n + 1]$。
这留给读者验证。
\end{proof}

\begin{lemma}
\label{hypercovering-lemma-homotopy}
设 $\mathcal{C}$ 是具有纤维积的位点，$X$ 是 $\mathcal{C}$ 的对象，
$K,L$ 是 $X$ 的超覆盖，$a,b : K \to L$ 是超覆盖的态射。
则存在超覆盖的态射 $c : K' \to K$，使得
$a \circ c$ 同伦于 $b \circ c$。
\end{lemma}

\begin{proof}
考虑交换图
$$
\xymatrix{
K' \ar@{=}[r]^-{def} \ar[rd]_c &
K \times_{(L \times L)} \Hom(\Delta[1], L)
\ar[r] \ar[d] & \Hom(\Delta[1], L) \ar[d] \\
& K \ar[r]^{(a, b)} & L \times L
}
$$
由 $\Hom(\Delta[1], L)$ 的函子性质，水平态射的复合对应于一个态射
$K' \times \Delta[1] \to L$；它定义 $c \circ a$ 与 $c \circ b$
之间的同伦。因此，只要证明 $K'$ 是 $X$ 的超覆盖，就可得到引理。
为此，把引理 \ref{hypercovering-lemma-funny-gamma} 应用于态射对
$K \to L \times L$ 和 $\Hom(\Delta[1], L) \to L \times L$。
引理 \ref{hypercovering-lemma-funny-gamma} 的条件 (1) 成立。引理
\ref{hypercovering-lemma-funny-gamma} 的条件 (2) 也成立，
因为 $\Hom(\Delta[1], L)_0 = L_1$；又因假设 $L$ 是超覆盖，态射
$(d^1_0, d^1_1) : L_1 \to L_0 \times L_0$ 是
$\text{SR}(\mathcal{C}, X)$ 中的覆盖。为证明引理
\ref{hypercovering-lemma-funny-gamma} 的条件 (3)，应用上述引理
\ref{hypercovering-lemma-hom-hypercovering}。根据该引理，引理
\ref{hypercovering-lemma-funny-gamma} 的条件 (3) 中的态射
$\gamma$ 是态射
$$
\Hom(\Delta[1] \times \Delta[n + 1], L)_0
\longrightarrow
\Hom(U, L)_0
$$
其中 $U \subset \Delta[1] \times \Delta[n + 1]$。由引理
\ref{hypercovering-lemma-add-simplices}，这是覆盖，故断言得证。
\end{proof}

\begin{remark}
\label{hypercovering-remark-contractible-category}
注意，证明的关键在于应用引理 \ref{hypercovering-lemma-add-simplices}。
这个引理完全是一般性的，并不在意单纯集的确切形状
（只要它们仅有有限多个非退化单形）。因此完全有理由期待如下结果：
给定任意态射 $a : K \times \partial \Delta[k] \to L$，其中 $K,L$
是超覆盖，则存在超覆盖的态射 $c : K' \to K$ 和态射
$g : K' \times \Delta[k] \to L$，使得
$g|_{K' \times \partial \Delta[k]} =
a \circ (c \times \text{id}_{\partial \Delta[k]})$.
换言之，超覆盖的范畴在适当意义下是可缩的。
\end{remark}

















\section{上同调与超覆盖}
\label{hypercovering-section-cohomology}

\noindent
设 $\mathcal{C}$ 是具有纤维积的位点，$X$ 是 $\mathcal{C}$ 的对象，
$\mathcal{F}$ 是 $\mathcal{C}$ 上的阿贝尔群层，$K,L$ 是 $X$ 的超覆盖。
若 $a,b : K \to L$ 是同伦映射，则
$\mathcal{F}(a), \mathcal{F}(b) : \mathcal{F}(K) \to \mathcal{F}(L)$
也是同伦映射，见《单纯方法》引理
\ref{simplicial-lemma-functorial-homotopy}。因此，它们在相伴上链复形的
上同调群上产生相同作用，见《单纯方法》引理
\ref{simplicial-lemma-homotopy-s-Q}。我们将用这一点定义遍历所有超覆盖的余极限。

\medskip\noindent
暂记 $\text{HC}(\mathcal{C}, X)$ 为如下范畴：其对象是 $X$ 的超覆盖，
其态射是 $X$ 的超覆盖之间的映射的同伦类。我们已经看到这是一个范畴，
而且不是“大”范畴，见引理 \ref{hypercovering-lemma-hypercoverings-set}。
$\text{HC}(\mathcal{C}, X)$ 的反范畴将作为下图的指标范畴；术语见
《范畴论》第 \ref{categories-section-limits} 节。考虑图
$$
\check{H}^i(-, \mathcal{F}) :
\text{HC}(\mathcal{C}, X)^{opp}
\longrightarrow
\textit{Ab}.
$$
由引理 \ref{hypercovering-lemma-product-hypercoverings}、\ref{hypercovering-lemma-homotopy}
以及上述关于同伦的说明，此图是有向的，见《范畴论》定义
\ref{categories-definition-directed}。因此余极限
$$
\check{H}^i_{\text{HC}}(X, \mathcal{F}) =
\colim_{K \in \text{HC}(\mathcal{C}, X)} \check{H}^i(K, \mathcal{F})
$$
有特别简单的描述（见所引位置）。

\begin{theorem}
\label{hypercovering-theorem-cohomology-hypercoverings}
设 $\mathcal{C}$ 是具有纤维积的位点，$X$ 是 $\mathcal{C}$ 的对象，
$i \geq 0$。函子
\begin{eqnarray*}
\textit{Ab}(\mathcal{C}) & \longrightarrow & \textit{Ab} \\
\mathcal{F} & \longmapsto & H^i(X, \mathcal{F}) \\
\mathcal{F} & \longmapsto & \check{H}^i_{\text{HC}}(X, \mathcal{F})
\end{eqnarray*}
典范同构。
\end{theorem}

\begin{proof}[使用谱序列的证明]
设对某个 $p \geq 0$ 有 $\xi \in H^p(X, \mathcal{F})$。
我们证明：对 $X$ 的某个超覆盖 $K$，$\xi$ 位于引理
\ref{hypercovering-lemma-cech-spectral-sequence} 的映射
$\check{H}^p(X, \mathcal{F}) \to H^p(X, \mathcal{F})$ 的像中。

\medskip\noindent
若 $p = 0$，这由引理 \ref{hypercovering-lemma-h0-cech} 成立。若 $p = 1$，
选取例 \ref{hypercovering-example-cech} 中 $X$ 的 {\v C}ech 超覆盖 $K$，
从位点 $\mathcal{C}$ 中满足 $\xi|_{U_i} = 0$ 的覆盖
$K_0 = \{U_i \to X\}$ 开始；见《位点上的上同调》引理
\ref{sites-cohomology-lemma-kill-cohomology-class-on-covering}。
此时立即由引理 \ref{hypercovering-lemma-cech-spectral-sequence} 的谱序列得出，
$\xi$ 来自 $\check{H}^1(K, \mathcal{F})$ 的一个元素。
一般地，选取 $X$ 的任意超覆盖 $K$，使得 $\xi$ 在
$\underline{H}^p(\mathcal{F})(K_0)$ 中的像为零（再次使用例
\ref{hypercovering-example-cech} 和《位点上的上同调》引理
\ref{sites-cohomology-lemma-kill-cohomology-class-on-covering}）。
由引理 \ref{hypercovering-lemma-cech-spectral-sequence} 的谱序列，$\xi$ 来自
$\check{H}^p(K, \mathcal{F})$ 的一个元素的障碍，是一列元素
$\xi_1, \ldots, \xi_{p - 1}$，其中
$\xi_q \in \check{H}^{p - q}(K, \underline{H}^q(\mathcal{F}))$
（更精确地说，是 $\xi_q$ 在这些群的某些子商中的像）。

\medskip\noindent
可以归纳地用细化替换超覆盖 $K$，使障碍
$\xi_1, \ldots, \xi_{p - 1}$ 的限制为零（而不仅是其在子商中的像为零，
故此处没有微妙之处）。事实上，假设已经达到
$\xi_{q + 1}, \ldots, \xi_{p - 1}$ 都为零的情形。注意，
$\xi_q \in \check{H}^{p - q}(K, \underline{H}^q(\mathcal{F}))$
是某个元素
$$
\tilde \xi_q \in
\underline{H}^q(\mathcal{F})(K_{p - q}) =
\prod H^q(U_i, \mathcal{F})
$$
的类，其中 $K_{p - q} = \{U_i \to X\}_{i \in I}$。设
$\xi_{q, i}$ 为 $\tilde \xi_q$ 在 $H^q(U_i, \mathcal{F})$ 中的分量。
由于 $q \geq 1$，可再次应用《位点上的上同调》引理
\ref{sites-cohomology-lemma-kill-cohomology-class-on-covering}，
选取位点的覆盖 $\{U_{i, j} \to U_i\}$，使每个限制
$\xi_{q, i}|_{U_{i, j}} = 0$。考虑范畴
$\text{SR}(\mathcal{C}, X)$ 的对象 $Z = \{U_{i, j} \to X\}$ 及其显然态射
$u : Z \to K_{p - q}$。显然 $u$ 是覆盖，见定义
\ref{hypercovering-definition-covering-SR}。由引理 \ref{hypercovering-lemma-covering}，存在 $X$
的超覆盖的态射 $L \to K$，使 $L_{p - q} \to K_{p - q}$ 经由 $u$
分解。于是显然，$\xi_q$ 在
$\underline{H}^q(\mathcal{F})(L_{p - q})$ 中的像为零。
由于引理 \ref{hypercovering-lemma-cech-spectral-sequence} 的谱序列是函子性的，
用 $L$ 替换 $K$ 后便达到 $\xi_q, \ldots, \xi_{p - 1}$ 均为零的情形。
如此继续，最终得到一个使它们全为零的超覆盖，因而 $\xi$ 位于映射
$\check{H}^p(X, \mathcal{F}) \to H^p(X, \mathcal{F})$ 的像中。

\medskip\noindent
设 $K$ 是 $X$ 的超覆盖，$\xi \in \check{H}^p(K, \mathcal{F})$，
且 $\xi$ 在引理 \ref{hypercovering-lemma-cech-spectral-sequence} 的映射
$\check{H}^p(X, \mathcal{F}) \to H^p(X, \mathcal{F})$ 下的像为零。
为完成定理证明，需说明存在超覆盖的态射 $L \to K$，使 $\xi$
在 $\check{H}^p(L, \mathcal{F})$ 中的限制为零。由引理
\ref{hypercovering-lemma-cech-spectral-sequence} 的谱序列，$\xi$ 在
$H^p(X, \mathcal{F})$ 中的像消失，意味着存在元素
$\xi_1, \ldots, \xi_{p - 2}$，其中
$\xi_q \in \check{H}^{p - 1 - q}(K, \underline{H}^q(\mathcal{F}))$
（更精确地说，是它们在某些子商中的像），使得谱序列中的像
$d_{q + 1}^{p - 1 - q, q}\xi_q$ 之和为 $\xi$。因此，完全按照上述机制，
可以找到超覆盖的态射 $L \to K$，使元素
$\xi_q$（$q = 1, \ldots, p - 2$）在
$\check{H}^{p - 1 - q}(L, \underline{H}^q(\mathcal{F}))$ 中的限制为零。
于是，由于态射 $L \to K$ 依照引理
\ref{hypercovering-lemma-cech-spectral-sequence} 诱导谱序列的态射，可得 $\xi$ 为零。
\end{proof}

\begin{proof}[不使用谱序列的证明]
我们已经看到 $i = 0$ 时的结论，见引理 \ref{hypercovering-lemma-h0-cech}。
已知函子 $H^i(X, -)$ 构成泛 $\delta$-函子，见《导出范畴》引理
\ref{derived-lemma-higher-derived-functors}。为证明定理，只需说明函子列
$\check{H}^i_{HC}(X, -)$ 构成 $\delta$-函子。事实上，已知
{\v C}ech 上同调在内射层上为零（引理
\ref{hypercovering-lemma-injective-trivial-cech}），于是可以应用《同调代数》引理
\ref{homology-lemma-efface-implies-universal}。

\medskip\noindent
设
$$
0 \to \mathcal{F} \to \mathcal{G} \to \mathcal{H} \to 0
$$
是 $\mathcal{C}$ 上阿贝尔层的短正合列。设
$\xi \in \check{H}^p_{HC}(X, \mathcal{H})$。选取 $X$ 的超覆盖 $K$
以及在上同调中代表 $\xi$ 的元素 $\sigma \in \mathcal{H}(K_p)$。
存在相应的复形正合列
$$
0 \to s(\mathcal{F}(K)) \to s(\mathcal{G}(K)) \to s(\mathcal{H}(K))
$$
但不能保证右端也有一个零；这正是妨碍我们直接应用蛇引理定义
$\delta(\xi)$ 的唯一因素。回忆，若 $K_p = \{U_i \to X\}$，则
$$
\mathcal{H}(K_p) = \prod \mathcal{H}(U_i)
$$
写 $\sigma =\prod \sigma_i$，其中 $\sigma_i \in \mathcal{H}(U_i)$。
由于 $\mathcal{G} \to \mathcal{H}$ 是层的满射，存在覆盖
$\{U_{i, j} \to U_i\}$，使 $\sigma_i|_{U_{i, j}}$ 是某个元素
$\tau_{i, j} \in \mathcal{G}(U_{i, j})$ 的像。考虑范畴
$\text{SR}(\mathcal{C}, X)$ 的对象 $Z = \{U_{i, j} \to X\}$ 及其显然态射
$u : Z \to K_p$。显然 $u$ 是覆盖，见定义
\ref{hypercovering-definition-covering-SR}。由引理 \ref{hypercovering-lemma-covering}，存在 $X$
的超覆盖的态射 $L \to K$，使 $L_p \to K_p$ 经由 $u$ 分解。
因此用 $L$ 替换 $K$ 后，可以假设 $\sigma$ 是某个元素
$\tau \in \mathcal{G}(K_p)$ 的像。注意 $d(\sigma) = 0$，
但不一定有 $d(\tau) = 0$。因而
$d(\tau) \in \mathcal{F}(K_{p + 1})$ 是闭链。在此情形下，定义
$\delta(\xi)$ 为闭链 $d(\tau)$ 在
$\check{H}^{p + 1}_{HC}(X, \mathcal{F})$ 中的类。

\medskip\noindent
此时有几件事需要验证：(a) $\delta(\xi)$ 不依赖于 $\tau$ 的选择；
(b) $\delta(\xi)$ 不依赖于使 $\sigma$ 可提升的超覆盖 $L \to K$
的选择；(c) $\delta(\xi)$ 不依赖于起初为表示 $\xi$ 所选的超覆盖和
$\sigma$。我们略去 (a)、(b)、(c) 的验证；其中对超覆盖选择的独立性，
实质上归结为引理 \ref{hypercovering-lemma-product-hypercoverings} 和
\ref{hypercovering-lemma-homotopy}。我们还略去验证 $\delta$ 关于 $\mathcal{C}$ 上
阿贝尔层短正合列之间的态射具有函子性。

\medskip\noindent
最后，必须验证在这个 $\delta$ 的定义下，上述阿贝尔层短正合列给出
{\v C}ech 上同调群的长正合列。先证明：若
$\delta(\xi) = 0$（$\xi$ 如上），则 $\xi$ 是某个元素
$\xi' \in \check{H}^p_{HC}(X, \mathcal{G})$ 的像。事实上，若
$\delta(\xi) = 0$，则沿用上述记号，$d(\tau)$ 的类在
$\check{H}^{p + 1}_{HC}(X, \mathcal{F})$ 中为零。因此存在超覆盖的态射
$L \to K$，使 $d(\tau)$ 限制所得的 $\mathcal{F}(L_{p + 1})$ 中的元素，
等于某个 $\upsilon \in \mathcal{F}(L_p)$ 的 $d(\upsilon)$。
这蕴含 $\tau|_{L_p} + \upsilon$ 构成闭链，并确定一个类
$\xi' \in \check{H}^p(L, \mathcal{G})$，它按要求映到 $\xi$。

\medskip\noindent
我们略去以下断言的证明：若
$\xi' \in \check{H}^{p + 1}_{HC}(X, \mathcal{F})$ 在
$\check{H}^{p + 1}_{HC}(X, \mathcal{G})$ 中映为零，则对某个
$\xi \in \check{H}^p_{HC}(X, \mathcal{H})$ 有
$\xi' = \delta(\xi)$。
\end{proof}

\noindent
下面用一个巧妙手法推出定理 \ref{hypercovering-theorem-cohomology-hypercoverings}
的 Verdier 情形。

\begin{proposition}
\label{hypercovering-proposition-cohomology-hypercoverings}
设 $\mathcal{C}$ 是具有纤维积和二元积的位点，
$\mathcal{F}$ 是 $\mathcal{C}$ 上的阿贝尔层，$i \geq 0$。则
\begin{enumerate}
\item 对每个 $\xi \in H^i(\mathcal{F})$，存在超覆盖 $K$，使 $\xi$
属于典范映射 $\check{H}^i(K, \mathcal{F}) \to H^i(\mathcal{F})$ 的像；
\item 若 $K,L$ 是超覆盖，且
$\xi_K \in \check{H}^i(K, \mathcal{F})$、
$\xi_L \in \check{H}^i(L, \mathcal{F})$ 映到
$H^i(\mathcal{F})$ 的同一元素，则存在超覆盖 $M$ 以及态射
$M \to K$ 和 $M \to L$，使 $\xi_K$ 与 $\xi_L$ 映到
$\check{H}^i(M, \mathcal{F})$ 的同一元素。
\end{enumerate}
换言之，忽略集合论问题，$\mathcal{F}$ 在 $\mathcal{C}$ 上的上同调群，
是 $\mathcal{F}$ 遍历所有超覆盖的 {\v C}ech 上同调群的余极限。
\end{proposition}

\begin{proof}
此结果是定理 \ref{hypercovering-theorem-cohomology-hypercoverings} 的直接推论。
事实上，可以人为地用一个稍大的位点 $\mathcal{C}'$ 替换
$\mathcal{C}$，使得 (I) $\mathcal{C}'$ 有终对象 $X$；(II)
$\mathcal{C}$ 中的超覆盖与 $\mathcal{C}'$ 中 $X$ 的超覆盖大致相同。
但由于问题本身的性质，需要处理相当多的记账细节。

\medskip\noindent
若 $\mathcal{C}$ 中具有固定靶的态射族 $\{U_i \to U\}$ 所对应的映射
$\coprod_{i \in I} h_{U_i} \to h_U$ 在层化后成为满射，则称它为
{\it 弱覆盖}。如下构造新位点 $\mathcal{C}'$：
\begin{enumerate}
\item 作为范畴，置
$\Ob(\mathcal{C}') = \Ob(\mathcal{C}) \amalg \{X\}$，
并从 $\mathcal{C}'$ 的每个对象向 $X$ 添加唯一态射；
\item 当 $\mathcal{C}$ 中存在纤维积和二元积时，
$\mathcal{C}'$ 具有相应的纤维积；
\item $\mathcal{C}'$ 的覆盖包括 $\mathcal{C}$ 的弱覆盖，以及满足如下条件的
$\{U_i \to X\}_{i \in I}$：或者某个 $i$ 有 $U_i = X$；或者对所有
$i$ 都有 $U_i \not = X$，且 $\mathcal{C}$ 上预层的映射
$\coprod h_{U_i} \to *$ 在 $\mathcal{C}$ 上层化后成为满射；
\item 应用《集合论》引理 \ref{sets-lemma-coverings-site} 限制覆盖，
从而得到位点 $\mathcal{C}'$。
\end{enumerate}
于是 $\Sh(\mathcal{C}') = \Sh(\mathcal{C})$，因为包含函子
$\mathcal{C} \to \mathcal{C}'$ 是特殊余连续函子
（见《位点与层》定义 \ref{sites-definition-special-cocontinuous-functor}）。
直接的验证从略。

\medskip\noindent
选取 $\mathcal{C}'$ 的覆盖 $\{U_i \to X\}$，使对所有 $i$，$U_i$ 都是
$\mathcal{C}$ 的对象（这是可能的，因为
$\mathcal{C} \to \mathcal{C}'$ 特殊余连续）。于是
$K_0 = \{U_i \to X\}$ 是上述位点 $\mathcal{C}'$ 中的覆盖。
把 $K_0$ 看成 $\text{SR}(\mathcal{C}', X)$ 的对象，并置
$K_{init} = \text{cosk}_0(K_0)$。于是 $K_{init}$ 是 $X$ 的超覆盖，
见例 \ref{hypercovering-example-cech}。注意，每个 $K_{init, n}$ 都具有
$\{W_j \to X\}$ 的形状，其中 $W_j \in \Ob(\mathcal{C})$。

\medskip\noindent
证明 (1)。选取
$\xi \in H^i(\mathcal{F}) = H^i(X, \mathcal{F}')$，其中
$\mathcal{F}'$ 是 $\mathcal{C}'$ 上与 $\mathcal{C}$ 上的
$\mathcal{F}$ 对应的阿贝尔层。由定理
\ref{hypercovering-theorem-cohomology-hypercoverings}，存在 $\mathcal{C}'$ 中 $X$
的超覆盖的态射 $K' \to K_{init}$，使 $\xi$ 来自
$\check{H}^i(K', \mathcal{F})$ 的一个元素。写
$K'_n = \{U_{n, j} \to X\}$。由于 $K'_n$ 映到
$K_{init, n}$，可见 $U_{n, j}$ 是 $\mathcal{C}$ 的对象。因此置
$K_n = \{U_{n, j}\}$，便可定义 $\text{SR}(\mathcal{C})$ 的单纯对象
$K$。由于 $\mathcal{C}'$ 中由 $\mathcal{C}$ 的态射族构成的覆盖，
是弱覆盖，可见 $K$ 是定义 \ref{hypercovering-definition-hypercovering-variant}
意义下的超覆盖。最后，由于 $\mathcal{F}'$ 是 $\mathcal{C}'$ 上
限制到 $\mathcal{C}$ 等于 $\mathcal{F}$ 的唯一层，{\v C}ech 复形
$s(\mathcal{F}(K))$ 与 $s(\mathcal{F}'(K'))$ 相同，于是 (1) 成立。
（略去与到上同调群的映射之相容性的验证。）

\medskip\noindent
证明 (2)。设 $K,L$ 是 $\mathcal{C}$ 中的超覆盖。令 $K',L'$ 为通过函子
$\text{SR}(\mathcal{C}) \to \text{SR}(\mathcal{C}', X)$,
$\{U_i\} \mapsto \{U_i \to X\}$ 由 $K,L$ 得到的
$\text{SR}(\mathcal{C}', X)$ 的单纯对象。与前面一样，{\v C}ech 复形相等，
于是得到 $\xi_{K'}$ 和 $\xi_{L'}$，它们映到 $\mathcal{C}'$ 上
$\mathcal{F}'$ 的同一上同调类。由于集合论问题，必要时扩大
$\mathcal{C}'$ 中覆盖的选择后，可以假设 $K',L'$ 是
$\mathcal{C}'$ 中 $X$ 的超覆盖；这由定义
\ref{hypercovering-definition-hypercovering-variant} 中的超覆盖定义，以及
$\mathcal{C}$ 中的弱覆盖给出 $\mathcal{C}'$ 中的覆盖这一事实成立。
由定理 \ref{hypercovering-theorem-cohomology-hypercoverings}，存在
$\mathcal{C}'$ 中 $X$ 的超覆盖 $M'$ 以及态射
$M' \to K'$、$M' \to L'$ 和 $M' \to K_{init}$，使
$\xi_{K'}$ 与 $\xi_{L'}$ 限制到
$\check{H}^i(M', \mathcal{F})$ 的同一元素。像上面一样展开此陈述，
便得 (2)。
\end{proof}





\section{空间的超覆盖}
\label{hypercovering-section-hypercoverings-spaces}

\noindent
即使对拓扑空间，上述理论也颇有趣。在这种情形下，可以具体写出超覆盖，
并看清结果实际表达的内容。

\medskip\noindent
设 $X$ 是拓扑空间。考虑《位点与层》例
\ref{sites-example-site-topological} 中的位点 $X_{Zar}$。
回忆，$X_{Zar}$ 的对象就是 $X$ 的开集，而 $X_{Zar}$ 的态射就是包含。
那么，位点 $X_{Zar}$ 上 $X$ 的超覆盖是什么？

\medskip\noindent
先展开定义 \ref{hypercovering-definition-SR}。$\text{SR}(X_{Zar}, X)$ 的一个对象，
就是一个集合 $I$ 以及对每个 $i \in I$ 给定的开集 $U_i \subset X$。
由于态射 $U_i \to X$ 不会混淆，把它记为 $\{U_i\}_{i \in I}$。
两个这样的对象之间的态射
$\{U_i\}_{i \in I} \to \{V_j\}_{j \in J}$，由集合映射
$\alpha : I \to J$ 给出，并满足对所有 $i \in I$ 有
$U_i \subset V_{\alpha(i)}$。这样的态射何时是覆盖？当且仅当对每个
$j \in J$ 都有
$V_j = \bigcup_{i\in I, \ \alpha(i) = j} U_i$（并且这是位点
$X_{Zar}$ 中的覆盖）。

\medskip\noindent
由以上内容得到位点 $X_{Zar}$ 中超覆盖的如下描述。
$X_{Zar}$ 中 $X$ 的一个超覆盖由以下数据给出：
\begin{enumerate}
\item 一个单纯集 $I$（见《单纯方法》第
\ref{simplicial-section-simplicial-set} 节）；
\item 对每个 $n \geq 0$ 和每个 $i \in I_n$，一个开集 $U_i \subset X$。
\end{enumerate}
用 $(I, \{U_i\})$ 表示这样一组数据。为使它成为 $X$ 的超覆盖，
要求满足以下性质：
\begin{itemize}
\item 对 $i \in I_n$ 和 $0 \leq a \leq n$，有
$U_i \subset U_{d^n_a(i)}$；
\item 对 $i \in I_n$ 和 $0 \leq a \leq n$，有
$U_i = U_{s^n_a(i)}$；
\item 有
\begin{equation}
\label{hypercovering-equation-covering-X}
X = \bigcup\nolimits_{i \in I_0} U_i,
\end{equation}
\item 对每个 $i_0, i_1 \in I_0$，有
\begin{equation}
\label{hypercovering-equation-covering-two}
U_{i_0} \cap U_{i_1} =
\bigcup\nolimits_{i \in I_1, \ d^1_0(i) = i_0, \ d^1_1(i) = i_1} U_i,
\end{equation}
\item 对每个 $n \geq 1$ 以及每个满足
$d^n_{b - 1}(i_a) = d^n_a(i_b)$（$0\leq a < b\leq n + 1$）的
$(i_0, \ldots, i_{n + 1}) \in (I_n)^{n + 2}$，有
\begin{equation}
\label{hypercovering-equation-covering-general}
U_{i_0} \cap \ldots \cap U_{i_{n + 1}} =
\bigcup\nolimits_{i \in I_{n + 1},
\ d^{n + 1}_a(i) = i_a, \ a = 0, \ldots, n + 1} U_i,
\end{equation}
\item 开覆盖 (\ref{hypercovering-equation-covering-X})、
(\ref{hypercovering-equation-covering-two}) 和 (\ref{hypercovering-equation-covering-general})
都属于 $\text{Cov}(X_{Zar})$（这是一个集合论条件，用来限制覆盖的指标集大小）。
\end{itemize}
例如，条件 (\ref{hypercovering-equation-covering-X}) 和
(\ref{hypercovering-equation-covering-two}) 应当可从《空间上的层》一章中认出，
而条件 (\ref{hypercovering-equation-covering-general}) 是其自然推广。

\begin{remark}
\label{hypercovering-remark-not-covering-set}
这个描述的一个特点是：若某个多重交
$U_{i_0} \cap \ldots \cap U_{i_{n + 1}}$ 为空，则右侧的覆盖可以是空覆盖。
因此，映射 $I_{n + 1} \to (\text{cosk}_n\text{sk}_n I)_{n + 1}$
不一定自动为满射。这意味着 $I$ 的几何实现可能是一个有趣的
（非可缩）空间。

\medskip\noindent
事实上，令 $I'_n \subset I_n$ 为由满足 $U_i \not = \emptyset$ 的单形
$i \in I_n$ 组成的子集。容易看出 $I' \subset I$ 是单纯子集，
且 $(I', \{U_i\})$ 是超覆盖。因此总能把一个超覆盖细化为所有开集
$U_i$ 都非空的超覆盖。
\end{remark}

\begin{remark}
\label{hypercovering-remark-repackage-into-simplicial-space}
再用另一种方式重新组织这些信息。设 $(I, \{U_i\})$ 是拓扑空间
$X$ 的超覆盖。由这些数据，置
$$
U_n = \coprod\nolimits_{i \in I_n} U_i,
$$
可构造单纯拓扑空间 $U_\bullet$；对给定的
$\varphi : [n] \to [m]$，令态射 $U(\varphi) : U_n \to U_m$
为来自各个 $i \in I_n$ 所对应包含 $U_i \subset U_{\varphi(i)}$ 的态射。
这个单纯拓扑空间带有增广 $\epsilon : U_\bullet \to X$。
凭借此态射，单纯空间 $U_\bullet$ 成为 $X$ 的超覆盖；沿此超覆盖，
存在 \cite[Expos\'e Vbis]{SGA4} 意义下的上同调下降。换言之，
$H^n(U_\bullet, \epsilon^*\mathcal{F}) = H^n(X, \mathcal{F})$。
（此处以后插入关于单纯空间上的上同调，以及用这种语言表述的上同调下降的参考。）
设 $\mathcal{F}$ 是 $X$ 上的阿贝尔层。此时，引理
\ref{hypercovering-lemma-cech-spectral-sequence} 的谱序列成为具有如下 $E_1$-项的谱序列：
$$
E_1^{p, q} = H^q(U_p, \epsilon_q^*\mathcal{F})
\Rightarrow
H^{p + q}(U_\bullet, \epsilon^*\mathcal{F}) = H^{p + q}(X, \mathcal{F})
$$
它把 $\epsilon^*\mathcal{F}$ 的全上同调与 $\mathcal{F}$ 在
$U_\bullet$ 各片上的上同调群作比较。（此处以后插入关于此谱序列的参考。）
\end{remark}

\noindent
在拓扑学中，我们常希望找到 $X$ 的超覆盖，使所有 $U_i$ 都来自
$X$ 的一个给定拓扑基，而且所有覆盖 (\ref{hypercovering-equation-covering-two})
和 (\ref{hypercovering-equation-covering-general}) 都来自一个给定的共终覆盖族。
下面给出两个示例性引理。

\begin{lemma}
\label{hypercovering-lemma-basis-hypercovering}
设 $X$ 是拓扑空间，$\mathcal{B}$ 是 $X$ 的拓扑基。
则存在 $X$ 的超覆盖 $(I, \{U_i\})$，使每个 $U_i$ 都属于
$\mathcal{B}$。
\end{lemma}

\begin{proof}
设 $n \geq 0$。所谓 $X$ 的一个{\it $n$-截断超覆盖}，是由一个
$n$-截断单纯集 $I$，以及对每个 $i \in I_a$、$0 \leq a \leq n$
给定的 $X$ 的开集 $U_i$ 所组成，并要求定义超覆盖的条件凡有意义时都成立。
换言之，只有当其中出现的所有单形都是满足 $a \leq n$ 的
$a$-单形时，才要求相应的包含关系和覆盖条件。若能证明：给定所有
$U_i \in \mathcal{B}$ 的 $n$-截断超覆盖 $(I, \{U_i\})$，可以在不添加
任何 $a \leq n$ 的 $a$-单形的情况下，把它扩张为
$(n + 1)$-截断超覆盖，那么引理即得证。具体如下。先考虑由
$I' = \text{sk}_{n + 1}(\text{cosk}_n I)$.
定义的 $(n + 1)$-截断单纯集 $I'$。回忆，
$$
I'_{n + 1} =
\left\{
\begin{matrix}
(i_0, \ldots, i_{n + 1}) \in (I_n)^{n + 2} \text{，使得}\\
d^n_{b - 1}(i_a) = d^n_a(i_b) \text{ 对所有 }0\leq a < b\leq n + 1
\end{matrix}
\right\}
$$
若 $i' \in I'_{n + 1}$ 是退化的，比如 $i' = s^n_a(i)$，则置
$U_{i'} = U_i$（无论如何，第二个条件已迫使我们这样做）。
此时还置 $J_{i'} = \{i'\}$。若 $i' \in I'_{n + 1}$ 是非退化的，
比如 $i' = (i_0, \ldots, i_{n + 1})$，则选取集合 $J_{i'}$
和开覆盖
\begin{equation}
\label{hypercovering-equation-choose-covering}
U_{i_0} \cap \ldots \cap U_{i_{n + 1}} =
\bigcup\nolimits_{i \in J_{i'}} U_i,
\end{equation}
其中对 $i \in J_{i'}$ 有 $U_i \in \mathcal{B}$。置
$$
I_{n + 1} = \coprod\nolimits_{i' \in I'_{n + 1}} J_{i'}
$$
存在典范映射 $\pi : I_{n + 1} \to I'_{n + 1}$；由构造，它在
$I'_{n + 1}$ 的退化单形集合上是双射。对 $i \in I_{n + 1}$，定义
$d^{n + 1}_a(i) = d^{n + 1}_a(\pi(i))$。对 $i \in I_n$，把
$s^n_a(i) \in I_{n + 1}$ 定义为位于退化单形
$s^n_a(i) \in I'_{n + 1}$ 上方的唯一单形。略去验证这定义了
$X$ 的一个 $(n + 1)$-截断超覆盖。
\end{proof}

\begin{lemma}
\label{hypercovering-lemma-quasi-separated-quasi-compact-hypercovering}
设 $X$ 是拓扑空间，$\mathcal{B}$ 是 $X$ 的拓扑基。假设
\begin{enumerate}
\item $X$ 是拟紧的；
\item 每个 $U \in \mathcal{B}$ 都是拟紧开集；
\item $X$ 中任意两个拟紧开集的交仍拟紧。
\end{enumerate}
则存在 $X$ 的超覆盖 $(I, \{U_i\})$，具有以下性质：
\begin{enumerate}
\item 每个 $U_i$ 都属于基 $\mathcal{B}$；
\item 每个 $I_n$ 都是有限集；特别地，
\item 覆盖 (\ref{hypercovering-equation-covering-X})、
(\ref{hypercovering-equation-covering-two}) 和 (\ref{hypercovering-equation-covering-general})
都是有限的。
\end{enumerate}
\end{lemma}

\begin{proof}
在引理 \ref{hypercovering-lemma-basis-hypercovering} 证明的构造中，若在
(\ref{hypercovering-equation-choose-covering}) 处选取由 $\mathcal{B}$ 中元素组成的
有限覆盖，便直接得到结论。细节从略。
\end{proof}







\section{构造超覆盖}
\label{hypercovering-section-hypercovering-sites}

\noindent
设 $\mathcal{C}$ 是位点。本节按如下方式看待
$\text{SR}(\mathcal{C})$ 的单纯对象。照例置 $K_n = K([n])$，
并用 $K(\varphi) : K_n \to K_m$ 表示与
$\varphi : [m] \to [n]$ 相伴的态射。可以写
$K_n = \{U_{n, i}\}_{i \in I_n}$。对
$\varphi : [m] \to [n]$，态射 $K(\varphi) : K_n \to K_m$
由映射 $\alpha(\varphi) : I_n \to I_m$ 以及对 $i \in I_n$ 的态射
$f_{\varphi, i} : U_{n, i} \to U_{m, \alpha(\varphi)(i)}$ 给出。
$K$ 是 $\text{SR}(\mathcal{C})$ 的单纯对象这一事实蕴含
$(I_n, \alpha(\varphi))$ 是单纯集，而且当
$\psi : [l] \to [m]$ 时，
$f_{\psi, \alpha(\varphi)(i)} \circ f_{\varphi, i} =
f_{\varphi \circ \psi, i}$。

\begin{lemma}
\label{hypercovering-lemma-split}
设 $\mathcal{C}$ 是位点，$K$ 是 $\text{SR}(\mathcal{C})$ 的
$r$-截断单纯对象。以下条件等价：
\begin{enumerate}
\item $K$ 是分裂的（《单纯方法》定义
\ref{simplicial-definition-split}）；
\item $f_{\varphi, i} : U_{n, i} \to U_{m, \alpha(\varphi)(i)}$
对 $r \geq n \geq 0$、满射 $\varphi : [m] \to [n]$ 和
$i \in I_n$ 是同构；
\item $f_{\sigma^n_j, i} : U_{n, i} \to U_{n + 1, \alpha(\sigma^n_j)(i)}$
对 $0 \leq j \leq n < r$ 和 $i \in I_n$ 是同构。
\end{enumerate}
对单纯对象也有同样结论，只需在 (2) 和 (3) 中置 $r = \infty$。
\end{lemma}

\begin{proof}
单纯集的分裂是唯一的，并由各次数 $n$ 中的非退化指标
$N(I_n)$ 给出，见《单纯方法》引理
\ref{simplicial-lemma-splitting-simplicial-sets}。
$\text{SR}(\mathcal{C})$ 的两个对象
$\{U_i\}_{i \in I}$ 和 $\{U_j\}_{j \in J}$ 的余积，按显然记号由
$\{U_l\}_{l \in I \amalg J}$ 给出。因此 $K$ 的分裂必由
$N(K_n) = \{U_i\}_{i \in N(I_n)}$ 给出。展开定义便得到 (1) 与 (2)
的等价性。任意满射 $\varphi : [m] \to [n]$ 都是态射
$\sigma^k_j$（$k = n, n + 1, \ldots, m - 1$）的复合，故 (2) 与 (3)
等价。
\end{proof}

\begin{lemma}
\label{hypercovering-lemma-hypercovering-object}
设 $\mathcal{C}$ 是具有纤维积的位点，
$\mathcal{B} \subset \Ob(\mathcal{C})$ 是子集。假设
\begin{enumerate}
\item $\mathcal{C}$ 的任意对象 $U$ 都有一个满足
$U_j \in \mathcal{B}$ 的覆盖 $\{U_j \to U\}_{j \in J}$；
\item 若 $\{U_j \to U\}_{j \in J}$ 是满足
$U_j \in \mathcal{B}$ 的覆盖，且 $\{U' \to U\}$ 是满足
$U' \in \mathcal{B}$ 的态射，则
$\{U_j \to U\}_{j \in J} \amalg \{U' \to U\}$ 是覆盖。
\end{enumerate}
则对 $\mathcal{C}$ 中的任意 $X$，存在 $X$ 的超覆盖 $K$，使得
$K_n = \{U_{n, i}\}_{i \in I_n}$，并且对所有 $i \in I_n$ 都有
$U_{n, i} \in \mathcal{B}$。
\end{lemma}

\begin{proof}
引理 \ref{hypercovering-lemma-basis-hypercovering} 的证明是本证明的预备练习；
建议读者先阅读那个证明。

\medskip\noindent
先用位点 $\mathcal{C}/X$ 替换 $\mathcal{C}$。这样做以后，可以假设
$X$ 是 $\mathcal{C}$ 的终对象，并且 $\mathcal{C}$ 具有所有有限极限
（《范畴论》引理 \ref{categories-lemma-finite-limits-exist}）。

\medskip\noindent
设 $n \geq 0$。所谓 $X$ 的一个
{\it $n$-截断 $\mathcal{B}$-超覆盖}，是
$\text{SR}(\mathcal{C})$ 的一个 $n$-截断单纯对象 $K$，满足：
对 $i \in I_a$、$0 \leq a \leq n$ 有
$U_{a, i} \in \mathcal{B}$；$K_0$ 是 $X$ 的覆盖；且
$K_{a + 1} \to (\text{cosk}_a \text{sk}_a K)_{a + 1}$
对 $a = 0, \ldots, n - 1$ 是定义
\ref{hypercovering-definition-covering-SR} 意义下的覆盖。

\medskip\noindent
由假设，$X$ 有一个满足 $U_i \in \mathcal{B}$ 的覆盖
$\{U_{0, i} \to X\}_{i \in I_0}$，故得到 $X$ 的一个
$0$-截断 $\mathcal{B}$-超覆盖。注意，$X$ 的任意
$0$-截断 $\mathcal{B}$-超覆盖都是分裂的，见引理 \ref{hypercovering-lemma-split}。

\medskip\noindent
若能证明：对 $n \geq 0$，给定 $X$ 的一个分裂
$n$-截断 $\mathcal{B}$-超覆盖 $K$，可把它扩张为 $X$ 的一个分裂
$(n + 1)$-截断 $\mathcal{B}$-超覆盖，则引理成立。

\medskip\noindent
构造该扩张。考虑 $\text{SR}(\mathcal{C})$ 的 $(n + 1)$-截断单纯对象
$K' = \text{sk}_{n + 1}(\text{cosk}_n K)$。写成
$$
K'_{n + 1} = \{U'_{n + 1, i}\}_{i \in I'_{n + 1}}
$$
由于 $K = \text{sk}_n K'$，对 $0 \leq a \leq n$ 有
$K_a = K'_a$。对每个 $i' \in I'_{n + 1}$，选取覆盖
\begin{equation}
\label{hypercovering-equation-choose-covering-B}
\{g_{n + 1, j} : U_{n + 1, j} \to U'_{n + 1, i'}\}_{j \in J_{i'}}
\end{equation}
使得对 $j \in J_{i'}$ 有 $U_{n + 1, j} \in \mathcal{B}$。
由引理中关于 $\mathcal{B}$ 的假设，这是可能的。对
$0 \leq m \leq n$，用 $N_m \subset I_m$ 表示非退化指标的子集。置
$$
I_{n + 1} =
\coprod\nolimits_{\varphi : [n + 1] \to [m]\text{ 满射，}0\leq m \leq n}
N_m \amalg
\coprod\nolimits_{i' \in I'_{n + 1}} J_{i'}
$$
对 $j \in I_{n + 1}$，置
$$
U_{n + 1, j} =
\left\{
\begin{matrix}
U_{m, i} & \text{若} &
j = (\varphi, i) & \text{其中} & \varphi : [n + 1] \to [m], i \in N_m \\
U_{n + 1, j} & \text{若} &
j \in J_{i'} & \text{其中} & i' \in I'_{n + 1}
\end{matrix}
\right.
$$
记号是显然的。置 $K_{n + 1} = \{U_{n + 1, j}\}_{j \in I_{n + 1}}$。
由构造，对所有 $j \in I_{n + 1}$，$U_{n + 1, j}$ 都属于
$\mathcal{B}$。定义相容映射
$$
I_{n + 1} \to I'_{n + 1}
\quad\text{和}\quad
K_{n + 1} \to K'_{n + 1}
$$
具体地，第一个映射由
$(\varphi, i) \mapsto \alpha'(\varphi)(i)$ 和
$(j \in J_{i'}) \mapsto i'$ 给出；第二个映射使用态射
$$
f'_{\varphi, i} : U_{m, i} \to U'_{n + 1, \alpha'(\varphi)(i)}
\quad\text{和}\quad
g_{n + 1, j} : U_{n + 1, j} \to U'_{n + 1, i'}
$$
我们断言，态射
$$
K_{n + 1} \to K'_{n + 1} =
(\text{cosk}_n \text{sk}_n K')_{n + 1} =
(\text{cosk}_n K)_{n + 1}
$$
是定义 \ref{hypercovering-definition-covering-SR} 意义下的覆盖。事实上，若
$i' \in I'_{n + 1}$，则或者 $i'$ 非退化，此时 $i'$ 在
$I_{n + 1}$ 中的逆像等于 $J_{i'}$，由选择
(\ref{hypercovering-equation-choose-covering-B}) 得到 $U'_{n + 1, i'}$ 的一个覆盖；
或者 $i'$ 退化，此时对唯一的偶对 $(\varphi, i)$，$i'$ 在
$I_{n + 1}$ 中的逆像为 $J_{i'} \amalg \{(\varphi, i)\}$，
由选择 (\ref{hypercovering-equation-choose-covering-B}) 和引理的假设 (2) 得到覆盖。

\medskip\noindent
为完成证明，还需定义与态射
$\varphi : [m] \to [n + 1]$、$0 \leq m \leq n$ 对应的态射
$K(\varphi) : K_{n + 1} \to K_m$，以及与态射
$\varphi : [n + 1] \to [m]$、$0 \leq m \leq n$ 对应的态射
$K(\varphi) : K_m \to K_{n + 1}$，并使它们满足适当的复合关系。
对第一类，使用复合
$$
K_{n + 1} \to K'_{n + 1} \xrightarrow{K'(\varphi)} K'_m = K_m
$$
来定义 $K(\varphi) : K_{n + 1} \to K_m$。对第二类，设给定
$\varphi : [n + 1] \to [m]$、$0 \leq m \leq n$。
如下定义相应态射 $K(\varphi) : K_m \to K_{n + 1}$：
\begin{enumerate}
\item 对 $i \in I_m$，存在唯一满射 $\psi : [m] \to [m_0]$
以及唯一的非退化 $i_0 \in I_{m_0}$，使得
$\alpha(\psi)(i_0) = i$\footnote{例如，若 $i$ 非退化，则
$m = m_0$ 且 $\psi = \text{id}_{[m]}$。}；
\item 置 $\varphi_0 = \psi_0 \circ \varphi : [n + 1] \to [m_0]$，
并把 $i \in I_m$ 映到 $(\varphi_0, i_0) \in I_{n + 1}$；换言之，
$\alpha(\varphi)(i) = (\varphi_0, i_0)$；
\item 态射
$f_{\varphi, i} : U_{m, i} \to U_{n + 1, \alpha(\varphi)(i)} = U_{m_0, i_0}$
是同构 $f_{\psi, i_0} : U_{m_0, i_0} \to U_{m, i}$ 的逆
（见引理 \ref{hypercovering-lemma-split}）。
\end{enumerate}
略去直接但繁琐的验证：这定义了 $X$ 的一个分裂
$(n + 1)$-截断 $\mathcal{B}$-超覆盖，扩张给定的 $n$-截断者。
事实上，除需验证对所有 $0 \leq a,b \leq n + 1$ 的
$\varphi : [a] \to [b]$，态射 $K(\varphi)$ 的复合正确之外，
其余一切均由以上构造显然成立。
\end{proof}

\begin{lemma}
\label{hypercovering-lemma-hypercovering-site}
设 $\mathcal{C}$ 是具有等化子和纤维积的位点，
$\mathcal{B} \subset \Ob(\mathcal{C})$ 是子集。假设
$\mathcal{C}$ 的任意对象都有一个各成员属于 $\mathcal{B}$ 的覆盖。
则存在超覆盖 $K$，使得 $K_n = \{U_i\}_{i \in I_n}$，
并且对所有 $i \in I_n$ 都有 $U_i \in \mathcal{B}$。
\end{lemma}

\begin{proof}
此证明与引理 \ref{hypercovering-lemma-hypercovering-object} 的证明几乎相同。
这里只说明不同之处。

\medskip\noindent
设 $n \geq 1$。所谓一个{\it $n$-截断 $\mathcal{B}$-超覆盖}，
是 $\text{SR}(\mathcal{C})$ 的一个 $n$-截断单纯对象 $K$，满足：
对 $i \in I_a$、$0 \leq a \leq n$ 有
$U_{a, i} \in \mathcal{B}$；并且
\begin{enumerate}
\item $F(K_0)^\# \to *$ 是满射；
\item $F(K_1)^\# \to F(K_0)^\# \times F(K_0)^\#$ 是满射；
\item $F(K_{a + 1})^\# \to F((\text{cosk}_a \text{sk}_a K)_{a + 1})^\#$
对 $a = 1, \ldots, n - 1$ 是满射。
\end{enumerate}
先显式构造一个分裂 $1$-截断 $\mathcal{B}$-超覆盖。

\medskip\noindent
取 $I_0 = \mathcal{B}$ 和 $K_0 = \{U\}_{U \in \mathcal{B}}$。
由关于 $\mathcal{B}$ 的假设，(1) 成立。置
$$
\Omega =
\{(U, V, W, a, b) \mid U, V, W \in \mathcal{B}, a : U \to V, b : U \to W\}
$$
再置 $I_1 = I_0 \amalg \Omega$。对 $i \in I_1$，若 $i \in I_0$，
置 $U_{1, i} = U_{0, i}$；若
$i = (U, V, W, a, b) \in \Omega$，置 $U_{1, i} = U$。
映射 $K(\sigma^0_0) : K_0 \to K_1$ 对应于包含
$\alpha(\sigma^0_0) : I_0 \to I_1$ 和对象上的恒等态射
$f_{\sigma^0_0, i} : U_{0, i} \to U_{1, i}$。映射
$K(\delta^1_0), K(\delta^1_1) : K_1 \to K_0$ 对应于两个映射
$I_1 \to I_0$：它们在 $I_0 \subset I_1$ 上为恒等，并分别把
$(U, V, W, a, b) \in \Omega \subset I_1$ 映到 $V$ 和 $W$。
相应态射 $f_{\delta^1_0, i}, f_{\delta^1_1, i} : U_{1, i} \to U_{0, i}$，
在 $i \in I_0$ 时为恒等；在
$i = (U, V, W, a, b) \in \Omega$ 时分别为 $a,b$。
(2) 成立的理由如下：$F(K_0)^\# \times F(K_0)^\#$ 在
$\mathcal{C}$ 的对象 $U$ 上的任意截面，在用某个覆盖的各成员替换
$U$ 后，都来自映射 $U \to F(K_0) \times F(K_0)$。
这又意味着存在 $V,W \in \mathcal{B}$ 和两个态射
$U \to V$、$U \to W$。进一步用某个覆盖的各成员替换 $U$ 后，
可以按要求假设 $U \in \mathcal{B}$。

\medskip\noindent
若能证明：对 $n \geq 1$，给定一个分裂
$n$-截断 $\mathcal{B}$-超覆盖 $K$，可以把它扩张为一个分裂
$(n + 1)$-截断 $\mathcal{B}$-超覆盖，则引理成立。这里的论证与引理
\ref{hypercovering-lemma-hypercovering-object} 的证明完全相同。除以下说明外，
精确细节从略。第一，在本引理的证明中不需要假设 (2)，因为无需态射
$K_{n + 1} \to (\text{cosk}_n K)_{n + 1}$ 本身是覆盖；只需它在相伴集合层上
诱导满射，而这由《位点与层》引理
\ref{sites-lemma-covering-surjective-after-sheafification} 得出。
第二，$\mathcal{C}$ 具有纤维积和等化子的假设保证
$\text{SR}(\mathcal{C})$ 具有纤维积和等化子，并且 $F$ 与它们交换
（引理 \ref{hypercovering-lemma-coprod-prod-SR}）。这足以保证所用的余骨架函子存在
（见《单纯方法》评注 \ref{simplicial-remark-existence-cosk} 和
《范畴论》引理 \ref{categories-lemma-fibre-products-equalizers-exist}）。
\end{proof}

\begin{lemma}
\label{hypercovering-lemma-hypercovering-morphism-sites}
设 $f : \mathcal{C} \to \mathcal{D}$ 是由函子
$u : \mathcal{D} \to \mathcal{C}$ 给出的位点态射。假设
$\mathcal{D}$ 和 $\mathcal{C}$ 具有等化子和纤维积，且 $u$ 与它们交换。
若 $\text{SR}(\mathcal{D})$ 的单纯对象 $K$ 是超覆盖，
则 $u(K)$ 是超覆盖。
\end{lemma}

\begin{proof}
如本节引言中那样写 $K_n = \{U_{n, i}\}_{i \in I_n}$，则 $u(K)$
是 $\text{SR}(\mathcal{C})$ 的对象，其中
$u(K_n) = \{u(U_i)\}_{i \in I_n}$。由《位点与层》引理
\ref{sites-lemma-pullback-representable-sheaf}，对
$U \in \Ob(\mathcal{D})$ 有 $f^{-1}h_U^\# = h_{u(U)}^\#$。
这意味着对所有 $n$，
$f^{-1}F(K_n)^\# = F(u(K_n))^\#$。验证定义
\ref{hypercovering-definition-hypercovering-variant} 中 $u(K)$ 成为超覆盖的条件
(1)、(2)、(3)。由于 $f^{-1}$ 是正合函子，可得
$$
F(u(K_0))^\# = f^{-1}F(K_0)^\# \to f^{-1}* = *
$$
作为满射的拉回是满射，故得到 (1)。类似地，
$$
F(u(K_1))^\# = f^{-1}F(K_1)^\# \to
f^{-1} (F(K_0) \times F(K_0))^\# = F(u(K_0))^\# \times F(u(K_0))^\#
$$
作为拉回是满射，故得到 (2)。对于条件 (3)，要用相同方法得出结论，
只需有
$$
F((\text{cosk}_n \text{sk}_n u(K))_{n + 1})^\# =
f^{-1}F((\text{cosk}_n \text{sk}_n K)_{n + 1})^\#
$$
以上说明 $f^{-1}F(-) = F(u(-))$。因此只需证明，对 $n \geq 1$，
$u$ 与定义 $(\text{cosk}_n \text{sk}_n K)_{n + 1}$ 时所用的极限交换。
由《单纯方法》评注 \ref{simplicial-remark-existence-cosk}，
这些极限是有限连通极限；由假设，$u$ 与它们交换。
\end{proof}

\begin{lemma}
\label{hypercovering-lemma-hypercovering-continuous-functor}
设 $\mathcal{C},\mathcal{D}$ 是位点，
$u : \mathcal{D} \to \mathcal{C}$ 是连续函子。假设
$\mathcal{D}$ 和 $\mathcal{C}$ 具有纤维积，且 $u$ 与纤维积交换。
设 $Y \in \mathcal{D}$，并设
$K \in \text{SR}(\mathcal{D}, Y)$ 是 $Y$ 的超覆盖。
则 $u(K)$ 是 $u(Y)$ 的超覆盖。
\end{lemma}

\begin{proof}
这比引理 \ref{hypercovering-lemma-hypercovering-morphism-sites} 的证明容易，
因为“是某个对象的超覆盖”这一概念更强，见定义
\ref{hypercovering-definition-hypercovering} 和 \ref{hypercovering-definition-covering-SR}。
事实上，由位点态射的定义，$u$ 把覆盖送到覆盖。只需验证，对
$n \geq 1$，$u$ 与定义
$(\text{cosk}_n \text{sk}_n K)_{n + 1}$ 时所用的极限交换。
这是显然的，因为诱导函子
$\mathcal{D}/Y \to \mathcal{C}/X$ 与所有有限极限交换
（而由《范畴论》引理 \ref{categories-lemma-finite-limits-exist}，
源和靶都具有所有有限极限）。
\end{proof}

\begin{lemma}
\label{hypercovering-lemma-w-contractible}
设 $\mathcal{C}$ 是位点，$\mathcal{B} \subset \Ob(\mathcal{C})$
是子集。假设
\begin{enumerate}
\item $\mathcal{C}$ 具有纤维积；
\item 对所有 $X \in \Ob(\mathcal{C})$，存在满足
$U_i \in \mathcal{B}$ 的有限覆盖 $\{U_i \to X\}_{i \in I}$；
\item 若 $\{U_i \to X\}_{i \in I}$ 是满足
$U_i \in \mathcal{B}$ 的有限覆盖，且 $U \to X$ 是满足
$U \in \mathcal{B}$ 的态射，则
$\{U_i \to X\}_{i \in I} \amalg \{U \to X\}$ 是覆盖。
\end{enumerate}
则对每个 $X$，存在 $X$ 的超覆盖 $K$，使每个
$K_n = \{U_{n, i} \to X\}_{i \in I_n}$ 中 $I_n$ 有限，且
$U_{n, i} \in \mathcal{B}$。
\end{lemma}

\begin{proof}
此引理是引理
\ref{hypercovering-lemma-quasi-separated-quasi-compact-hypercovering} 的位点版本。
证明时完全依照引理 \ref{hypercovering-lemma-hypercovering-object} 的证明，
并注意以下两点：
\begin{enumerate}
\item[(a)] 选取初始覆盖 $\{U_{0, i} \to X\}_{i \in I_0}$，
使 $U_{0, i} \in \mathcal{B}$ 且指标集 $I_0$ 有限；
\item[(b)] 选取覆盖 (\ref{hypercovering-equation-choose-covering-B}) 时，
取 $J_{i'}$ 为有限集。
\end{enumerate}
读者容易看出，经过这些修改，最终对所有 $n$ 都得到有限指标集 $I_n$。
\end{proof}

\begin{remark}
\label{hypercovering-remark-taking-disjoint-unions}
设 $\mathcal{C}$ 是位点，$K,L$ 是 $\text{SR}(\mathcal{C})$ 的对象。
写 $K = \{U_i\}_{i \in I}$、$L = \{V_j\}_{j \in J}$。
假设 $U = \coprod_{i \in I} U_i$ 和
$V = \coprod_{j \in J} V_j$ 存在。于是得到
$$
\Mor_{\text{SR}(\mathcal{C})}(K, L) \longrightarrow \Mor_\mathcal{C}(U, V)
$$
如下。给定由 $\alpha : I \to J$ 和
$f_i : U_i \to V_{\alpha(i)}$ 给出的 $f : K \to L$，得到函子的变换
$$
\Mor_\mathcal{C}(V, -) =
\prod\nolimits_{j \in J} \Mor_\mathcal{C}(V_j, -)
\to
\prod\nolimits_{i \in I} \Mor_\mathcal{C}(U_i, -) =
\Mor_\mathcal{C}(U, -)
$$
它把 $(g_j)_{j \in J}$ 送到
$(g_{\alpha(i)} \circ f_i)_{i \in I}$。因此 Yoneda 引理给出相应映射
$U \to V$。当然，$U \to V$ 通过态射 $f_i$ 把直和项 $U_i$
映入直和项 $V_{\alpha(i)}$。
\end{remark}

\begin{remark}
\label{hypercovering-remark-take-unions-hypercovering}
设 $\mathcal{C}$ 是具有纤维积和等化子的位点，$K$ 是超覆盖。
写 $K_n = \{U_{n, i}\}_{i \in I_n}$。假设
\begin{enumerate}
\item[(a)] $U_n = \coprod_{i \in I_n} U_{n, i}$ 存在；
\item[(b)] $\coprod_{i \in I_n} h_{U_{n, i}} \to h_{U_n}$
在层化后诱导同构。
\end{enumerate}
于是得到 $\text{SR}(\mathcal{C})$ 的另一个单纯对象 $L$，其中
$L_n = \{U_n\}$，见评注 \ref{hypercovering-remark-taking-disjoint-unions}。
我们断言 $L$ 是超覆盖。为此验证定义
\ref{hypercovering-definition-hypercovering-variant} 的条件 (1)、(2)、(3)。
条件 (1) 由 (b) 和 $K$ 的条件 (1) 得出；条件 (2) 完全类似。
条件 (3) 由下式得出：
\begin{align*}
F((\text{cosk}_n \text{sk}_n L)_{n + 1})^\#
& =
((\text{cosk}_n \text{sk}_n F(L)^\#)_{n + 1}) \\
& =
((\text{cosk}_n \text{sk}_n F(K)^\#)_{n + 1}) \\
& =
F((\text{cosk}_n \text{sk}_n K)_{n + 1})^\#
\end{align*}
其中 $n \geq 1$；因此完全如 (1)、(2) 那样，$K$ 的相应条件蕴含
$L$ 的相应条件。注意，$F$ 与连通极限交换且层化正合，这证明第一和
最后一个等式；中间的等式由 (b) 中
$F(K)^\# = F(L)^\#$ 得出。
\end{remark}

\begin{remark}
\label{hypercovering-remark-take-unions-hypercovering-X}
设 $\mathcal{C}$ 是位点，$X \in \Ob(\mathcal{C})$。假设
$\mathcal{C}$ 具有纤维积，并设 $K$ 是 $X$ 的超覆盖。
写 $K_n = \{U_{n, i}\}_{i \in I_n}$。假设
\begin{enumerate}
\item[(a)] $U_n = \coprod_{i \in I_n} U_{n, i}$ 存在；
\item[(b)] 给定 $\text{SR}(\mathcal{C})$ 中的态射
$(\alpha, f_i) : \{U_i\}_{i \in I} \to \{V_j\}_{j \in J}$ 和
$(\beta, g_k) : \{W_k\}_{k \in K} \to \{V_j\}_{j \in J}$
，且
$U = \coprod U_i$、$V = \coprod V_j$ 和 $W = \coprod W_j$
存在，则 $U \times_V W =
\coprod_{(i, j, k), \alpha(i) = j = \beta(k)} U_i \times_{V_j} W_k$；
\item[(c)] 若 $(\alpha, f_i) : \{U_i\}_{i \in I} \to \{V_j\}_{j \in J}$
是定义 \ref{hypercovering-definition-covering-SR} 意义下的覆盖，且
$U = \coprod U_i$ 和 $V = \coprod V_j$ 存在，则评注
\ref{hypercovering-remark-taking-disjoint-unions} 中相应的态射 $U \to V$
是 $\mathcal{C}$ 的覆盖。
\end{enumerate}
于是得到 $\text{SR}(\mathcal{C})$ 的另一个单纯对象 $L$，其中
$L_n = \{U_n\}$，见评注 \ref{hypercovering-remark-taking-disjoint-unions}。
我们断言 $L$ 是 $X$ 的超覆盖。为此验证定义
\ref{hypercovering-definition-hypercovering} 的条件 (1)、(2)。条件 (1) 由 (c)
和 $K$ 的条件 (1) 得出，因为 $K$ 的条件 (1) 表明
$K_0 = \{U_{0, i}\}_{i \in I_0}$ 是定义
\ref{hypercovering-definition-covering-SR} 意义下 $\{X\}$ 的覆盖。
条件 (2) 成立，是因为 $\mathcal{C}/X$ 具有所有有限极限，故
$\text{SR}(\mathcal{C}/X)$ 也具有所有有限极限，而条件 (b) 表明
“取不交并”的构造与这些有限极限交换。因此态射
$$
L_{n + 1} \longrightarrow (\text{cosk}_n \text{sk}_n L)_{n + 1}
$$
是覆盖，因为它是把“取不交并”函子应用于态射
$$
K_{n + 1} \longrightarrow (\text{cosk}_n \text{sk}_n K)_{n + 1}
$$
所得的结果；由 $K$ 的条件 (2)，后者被假设为定义
\ref{hypercovering-definition-covering-SR} 意义下的覆盖。这一说法有意义，
因为性质 (b) 特别保证：若从 $X$ 上一张可取不交并的半可表对象有限图出发，
则该图在 $\text{SR}(\mathcal{C}/X)$ 中的极限仍是 $X$ 上可取不交并的
半可表对象。
\end{remark}
